% =====================================================================
%  第 9 章 临界水平 \hat g-模的结构（续）
%  §9.4 临界水平 \hat g-模的表示范畴（印刷 p.274 下–p.281）
%  §9.5 Verma 模的自同态（印刷 p.282–p.289）
%  §9.6 Weyl 模的自同态（印刷 p.289–p.304）
%  对应 PDF p.284–p.314。本文件不含 \chapter。
%
%  【接缝说明】PDF p.284 上半页仍属 §9.3（“Miura opers with regular singularities”）：
%  含 Springer 纤维双射的收尾、定理 9.3.7、式 (9.3-18) 及其后解释。U12 止于 PDF p.283，
%  这段续文放在第一个 section 标题之前；经 Lead 核对 ch09-1 不含这些内容，**不得删除**。
% =====================================================================

与同构 (9.3-17) 的情形一样，这只是 $\mathbb{C}$-点层面的一一对应。纤维 $\mu^{-1}(\chi)$ 是各层之并：
\[
    \mu^{-1}(\chi)_w=\mu^{-1}(\chi)\cap\operatorname{Conn}(\Omega^{\check\rho})^{\check\rho-w(\check\lambda+\check\rho)}_D,
    \qquad w\in W,
\]
而每一层作为簇同构于子簇
\[
    \operatorname{Sp}^w_{\operatorname{Res}_{\check\lambda}(\chi)}
    =\operatorname{Sp}_{\operatorname{Res}_{\check\lambda}(\chi)}\cap S_w,
\]
后者由那些与 $\mathfrak{b}$ 处于相对位置 $w$ 的 Borel 子代数组成；这里 $S_w$ 是 $G/B$ 中对应于 $w$ 的 Schubert 胞腔（见引理 9.3.2 的证明）。

\begin{theorem}
    对任意 $\chi\in\Op^{\mathrm{nilp},\check\lambda}_G$，Miura 变换 $\mu$ 在 $\chi$ 上的纤维 $\mu^{-1}(\chi)$ 是各层之并
    \[
        \mu^{-1}(\chi)_w\subset\operatorname{Conn}(\Omega^{\check\rho})^{\check\rho-w(\check\lambda+\check\rho)}_D,
        \qquad w\in W.
    \]
    层 $\mu^{-1}(\chi)_w$ 同构于
    \[
        \operatorname{Sp}^w_{\operatorname{Res}_{\check\lambda}(\chi)}
        =\operatorname{Sp}_{\operatorname{Res}_{\check\lambda}(\chi)}\cap S_w,
    \]
    其中 $\operatorname{Sp}_{\operatorname{Res}_{\check\lambda}(\chi)}$ 是 $\operatorname{Res}_{\check\lambda}(\chi)$ 的 Springer 纤维。因此 $\mu^{-1}(\chi)$ 的点集与 $\operatorname{Sp}_{\operatorname{Res}_{\check\lambda}(\chi)}$ 的点集一一对应。
\end{theorem}

例如若 $\operatorname{Res}_{\check\lambda}(\chi)=0$，则 $\operatorname{Sp}_{\operatorname{Res}_{\check\lambda}(\chi)}$ 是旗簇 $G/B$，而 $\operatorname{Sp}^w_{\operatorname{Res}_{\check\lambda}(\chi)}$ 是 Schubert 胞腔 $S_w$。于是此时 $\mu^{-1}(\chi)$ 与 $G/B$ 的点集一一对应。

注意 $\tilde{\mathfrak{n}}_{w_0}$ 是单点（即 $\mathfrak{b}'=\mathfrak{b}$），故 $\mu^{-1}(\chi)_{w_0}=\operatorname{Sp}^{w_0}_{\operatorname{Res}_{\check\lambda}(\chi)}$ 也由单点组成。因此有同构
\begin{equation}
    \operatorname{Conn}(\Omega^{\check\rho})^{\check\rho-w_0(\check\lambda+\check\rho)}_D
    \simeq\Op^{\mathrm{nilp},\check\lambda}_G .
    \tag{9.3-18}
\end{equation}
这其实是命题 9.3.1 陈述的特例（另见推论 9.2.5），因为若 $\check\lambda$ 是整优余权，则余权 $\check\mu=w_0(\check\lambda+\check\rho)-\check\rho$ 对一切 $\alpha\in\Delta_+$ 满足 $\langle\alpha,\check\mu\rangle\notin\Z_+$。

% [Tutor Note: p.274 的接续段中“带 check 的记号”必须与 §9.4 起“不带 check 的记号”分开看。此处仍在 G-侧（$H$-丛 $\Omega^{\check\rho}$、余权 $\check\lambda,\check\rho$）；从 §9.4 起作者换到 ${}^LG$-侧（${}^LH$-丛 $\Omega^\rho$、权 $\lambda,\rho$）。原书个别位置混用了两套记法（如 p.291 的 $\Op^{\mathrm{nilp},\lambda}$ 与 p.300 的 $\Op^{\check\lambda}$），译文一律照原页符号，不擅自统一。]

\section{临界水平 $\gghat$-模的表示范畴}

\begin{keypoints}
    \item \textbf{核心思想：}临界水平的光滑 $\gghat_{\kappa_c}$-模范畴“生活”在 $\Op_{{}^LG}(D^\times)$ 之上；本节把若干子范畴对应到 $\Op^{\mathrm{RS}}_{{}^LG}(D)$ 的（形式）子簇上，并给出 Wakimoto 模所在的范畴。
    \item \textbf{数学工具：}中心的 $\Z$-分次（由 $L_0=-t\partial_t$ 给出）；Miura 变换；有限维 Harish--Chandra 同态；半无限抛物诱导得到的 Wakimoto 模；Springer 纤维（斜变子空间）。
    \item \textbf{关键结论：}命题 9.4.2 给出 $\operatorname{Ind}(M)$ 的概形论支撑；命题 9.4.1 与式 (9.4-3) 把中心的最低次部分识别为 $\operatorname{Fun}\Op^{\mathrm{RS}}_{{}^LG}(D)$；Wakimoto 模 $W_{\overline{\nabla}}$ 给出以 Springer 纤维为参数族的对象。
    \item \textbf{记号提醒：}本节起 $G$ 换成 Langlands 对偶群 ${}^LG$（取伴随型）；$\Op^{\mathrm{RS}}_{{}^LG}(D)_\varpi$ 的下标 $\varpi$ 是留数，与 9.3 节的 $\check\lambda$ 不同。
\end{keypoints}

前几节定义了去心圆盘 $D^\times$ 上 $G$-oper 空间 $\Op_G(D^\times)$ 的某些子簇。现在把 $G$ 换成 Langlands 对偶群 ${}^LG$（如前，假定它取伴随型）。特别地，这意味着我们从 $H$-丛 $\Omega^{\check\rho}$ 上联络的空间 $\operatorname{Conn}(\Omega^{\check\rho})_{D^\times}$ 及其上述各子空间，换到 ${}^LH$-丛 $\Omega^\rho$ 上联络的空间 $\operatorname{Conn}(\Omega^\rho)$ 及相应的子空间。

由定理 8.3.1，空间 $\Op_{{}^LG}(D^\times)$ 同构于临界水平完备泛包络代数 $\gghat$ 的中心 $Z(\gghat)$ 的谱。于是，全体光滑 $\gghat_{\kappa_c}$-模所成的范畴“生活”在 $\Op_{{}^LG}(D^\times)$ 之上。本节将识别该范畴的某些子范畴，它们“生活”在子空间 $\Op^{\mathrm{RS}}_{{}^LG}(D)_\varpi\subset\Op_{{}^LG}(D^\times)$ 之上。\transnote{「生活」在此有确定含义：由定理 8.3.1，$Z(\gghat)\simeq\operatorname{Fun}\Op_{{}^LG}(D^\times)$，故一个光滑 $\gghat_{\kappa_c}$-模上的中心作用可通过某个子簇上的函数代数分解，该子簇就是它的支撑。本节把诱导模收到 $\Op^{\mathrm{RS}}_{{}^LG}(D)$，§9.5 把 Verma 模收到单层 $\Op^{\mathrm{RS}}_{{}^LG}(D)_{\varpi(-\nu-\rho)}$，§9.6 把 Weyl 模再收到 $\Op^{\check\lambda}_{{}^LG}$；三条结论形状相同，只是子簇逐次变窄。落点：§9.4 开篇，原书 p.274–275。[来源：本书 p.274–275；定理 8.3.1。]}

\subsection{与有限维 Harish--Chandra 同态的相容性}

先描述 Miura 变换与 Harish--Chandra 同构 $Z(\mathfrak{g})\simeq(\operatorname{Fun}\mathfrak{h}^*)^W$ 之间的某种相容性，其中 $Z(\mathfrak{g})$ 是 $U(\mathfrak{g})$ 的中心。

考虑 $\widetilde{U}_{\kappa_c}(\gghat)$ 的中心 $Z(\gghat)$。算子 $L_0=-t\partial_t\in\operatorname{Der}\mathcal{O}$ 给出 $Z(\gghat)$ 上一个自然的 $\Z$-分次。

我们把任一 $\Z$-分次对象的 0 次部分记以下标 0，把所有负次元素张成的子空间记以下标 $<0$。

接着考虑 Wakimoto 模 $W_{\chi(t)}$，其中
\[
    \chi(t)=\frac{\lambda-\rho}{t},\qquad \lambda\in\mathfrak{h}^*
\]
（见定理 6.4.1 的证明）。这些 Wakimoto 模是 $\Z$-分次的，其 0 次部分 $(W_{\chi(t)})_0$ 是子空间 $\mathbb{C}[a^*_{\alpha,0}|0\rangle]_{\alpha\in\Delta_+}\subset M_{\mathfrak{g}}\subset W_{\chi(t)}$。李代数 $\mathfrak{g}$ 保持该子空间；由 Wakimoto 模的构造立即可知，作为 $\mathfrak{g}$-模，$(W_{\chi(t)})_0$ 同构于反变 Verma 模 $M^*_{\lambda-\rho}$。代数
\[
    Z'=Z(\gghat)_0/(Z(\gghat)\cdot Z(\gghat)_{<0})_0
\]
自然地作用在 $(W_{\chi(t)})_0$ 上并与 $\mathfrak{g}$ 交换。又因它的作用必然经由 $U(\mathfrak{g})$ 的作用分解，故经由 $Z(\mathfrak{g})$ 分解。因此让 $\lambda\in\mathfrak{h}^*$ 变动，便得到同态
\begin{equation}
    Z'\to(\operatorname{Fun}\mathfrak{h}^*)^W,
    \tag{9.4-1}
\end{equation}
它经由 Harish--Chandra 同态 $Z(\mathfrak{g})\to(\operatorname{Fun}\mathfrak{h}^*)^W$ 分解。

另一方面，回忆 §9.1.1 中的空间
\[
    \Op^{\mathrm{RS}}_{{}^LG}(D)\subset\Op_{{}^LG}(D^\times),
\]
即具有正则奇点的 ${}^LG$-oper 空间，及其由留数 $\varpi\in{}^L\mathfrak{h}/W=\mathfrak{h}^*/W$ 确定的子空间 $\Op^{\mathrm{RS}}_{{}^LG}(D)_\varpi$（回忆 $\mathfrak{h}^*={}^L\mathfrak{h}$）。\transnote{$\Op^{\mathrm{RS}}_{{}^LG}(D)_\varpi$ 是留数映射 $\operatorname{res}:\Op^{\mathrm{RS}}_{{}^LG}(D)\to\mathfrak{h}^*/W$ 的纤维，$\varpi$ 取遍 $\mathfrak{h}^*/W$；因 $\varpi$ 是 $W$-轨道，写 $\varpi\in{}^L\mathfrak{h}/W=\mathfrak{h}^*/W$ 与写 $\mathfrak{h}^*/W$ 中的点是一回事。每一层 $\Op^{\mathrm{RS}}_{{}^LG}(D)_\varpi$ 仍是无限维的（高阶系数自由），但它的 0 次函数代数只有限维，正是 (9.4-2) 右端的 $(\operatorname{Fun}\Op^{\mathrm{RS}}_{{}^LG}(D))_0\simeq(\operatorname{Fun}\mathfrak{h}^*)^W$。落点：§9.4.1，原书 p.275–276。[来源：本书 p.275–276 与 §9.1.1 的定义。]}

% [Tutor Note: p.275 的 $M^*_{\lambda-\rho}$ 与第 8 章的 $b_{i,n}\mapsto-u_{i,n}$ 符号约定要一起读：Wakimoto 模 $W_{\chi(t)}$ 的 0 次部分实现的是 $M^*_{\lambda-\rho}$，因此这里的 $\lambda$ 与 §9.5 中 Verma 模 $\mathbb{M}_{w(\lambda+\rho)-\rho}$ 的 $\lambda$ 相差 $\rho$ 平移。核对中心特征标时最容易在这一步丢掉 $\rho$。]

由算子 $L_0=-t\partial_t\in\operatorname{Der}\mathcal{O}$ 诱导出 $\operatorname{Fun}\Op_{{}^LG}(D^\times)$ 上的 $\Z$-分次。若把 oper 联络写成引理 4.2.2 的规范形式
\[
    \nabla=\partial_t+p_{-1}+\sum_{j=1}^{\ell}v_j(t)\cdot p_j,
    \qquad
    v_j(t)=\sum_{n\in\Z}v_{j,n}t^{-n-1},
\]
则由 §4.2.4 的变换公式得到
\[
    L_0\cdot v_{j,n}=(-n+d_j)v_{j,n}.
\]
因此生成元 $v_{j,n}$ 都是齐次的，严格负次的生成元是 $v_{j,n}$（$n>d_j$）；由推论 9.1.2，它们生成 $\operatorname{Fun}\Op_G(D^\times)$ 中 $\Op^{\mathrm{RS}}_G(D)$ 的理想。故 $\Op^{\mathrm{RS}}_G(D)$ 的理想由 $\operatorname{Fun}\Op_{{}^LG}(D^\times)_{<0}$ 生成。显然留数映射
\[
    \operatorname{res}:\Op^{\mathrm{RS}}_G(D)\to\mathfrak{h}^*/W
\]
给出同构
\[
    (\operatorname{Fun}\mathfrak{h}^*)^W\simeq(\operatorname{Fun}\Op^{\mathrm{RS}}_{{}^LG}(D))_0 .
\]

我们还有空间
\[
    \operatorname{Conn}(\Omega^\rho)^{\mathrm{RS}}_D\subset\operatorname{Conn}(\Omega^\rho)_{D^\times},
\]
即 §9.3.1 中定义在 $D^\times$ 上 ${}^LH$-丛 $\Omega^\rho$ 的具有正则奇点的联络空间。Miura 变换限制为映射（见 (9.3-4)）
\[
    \mu^{\mathrm{RS}}:\operatorname{Conn}(\Omega^\rho)^{\mathrm{RS}}_D\to\Op^{\mathrm{RS}}_{{}^LG}(D).
\]
算子 $L_0=-t\partial_t\in\operatorname{Der}\mathcal{O}$ 诱导出 $\operatorname{Fun}\operatorname{Conn}(\Omega^\rho)_{D^\times}$ 上的 $\Z$-分次。若把 $\overline{\nabla}\in\operatorname{Conn}(\Omega^\rho)_{D^\times}$ 写成 $\overline{\nabla}=\partial_t+\mathbf{u}(t)$，其中
\[
    u_i(t)=\alpha_i(\mathbf{u}(t))=\sum_{n<0}u_{i,n}t^{-n-1},
\]
则
\[
    L_0\cdot u_{i,n}=-nu_{i,n}.
\]
因此生成元 $u_{i,n}$ 都是齐次的，严格负次的是 $v_{i,n}$（$n>0$）。\transnote{原文此处两处记号照页图保留：一是「严格负次的是 $v_{i,n}$」中的 $v_{i,n}$ 应为 $u_{i,n}$（同段上一句已写明生成元是 $u_{i,n}$，且 $L_0u_{i,n}=-nu_{i,n}$ 表明严格负次者当是 $n>0$ 的 $u_{i,n}$）；二是本页把 oper 的下标写作 $G$（$\Op^{\mathrm{RS}}_G(D)$、$\operatorname{res}_{{}^L\mathfrak{h}}:\Op^{\mathrm{RS}}_G(D)\to\mathfrak{h}^*$），而此处已换到对偶群一侧，应读作 ${}^LG$（同页后面出现的同构即写 $\operatorname{Fun}\Op^{\mathrm{RS}}_{{}^LG}(D)$）。[来源：本书 p.276；并见 p.275 的记号提醒。]}它们生成 $\operatorname{Fun}\operatorname{Conn}(\Omega^\rho)_{D^\times}$ 中 $\operatorname{Conn}(\Omega^\rho)^{\mathrm{RS}}_D$ 的理想。§9.3.1 中定义的留数映射
\[
    \operatorname{res}_{{}^L\mathfrak{h}}:\Op^{\mathrm{RS}}_G(D)\to\mathfrak{h}^*
\]
给出同构
\[
    \operatorname{Fun}\mathfrak{h}^*\to(\operatorname{Fun}\operatorname{Conn}(\Omega^\rho)^{\mathrm{RS}}_D)_0 .
\]

% [Tutor Note: p.276 有两处照原文保留的瑕疵，都在记号层面。(1)“严格负次的生成元是 $v_{i,n}$（$n>0$）”中的 $v_{i,n}$ 应为 $u_{i,n}$（同段上一句已说生成元是 $u_{i,n}$）。(2) 本页把下标写作 $G$（$\Op^{\mathrm{RS}}_G(D)$、$\operatorname{res}_{{}^L\mathfrak{h}}:\Op^{\mathrm{RS}}_G(D)\to\mathfrak{h}^*$），而此处已在 ${}^LG$-侧，$G$ 疑为笔误，应读作 ${}^LG$。译文照原页，读者按上下文理解。]

我们有交换图 (9.3-5)
\begin{equation*}
    \begin{CD}
        \operatorname{Conn}(\Omega^\rho)^{\mathrm{RS}}_D @>\mu^{\mathrm{RS}}>>
            \Op^{\mathrm{RS}}_{{}^LG}(D)\\
        @V{\operatorname{res}_{{}^L\mathfrak{h}}}VV @VV{\operatorname{res}}V\\
        \mathfrak{h}^* @>>> \mathfrak{h}^*/W
    \end{CD}
\end{equation*}
它给出代数同态的交换图
\begin{equation}
    \begin{CD}
        \operatorname{Fun}\mathfrak{h}^* @>>>
            (\operatorname{Fun}\operatorname{Conn}(\Omega^\rho)^{\mathrm{RS}}_D)_0\\
        @AAA @AAA\\
        (\operatorname{Fun}\mathfrak{h}^*)^W @>>>
            (\operatorname{Fun}\Op^{\mathrm{RS}}_{{}^LG}(D))_0
    \end{CD}
    \tag{9.4-2}
\end{equation}
现在，由于 $\Op^{\mathrm{RS}}_{{}^LG}(D)$ 在 $\operatorname{Fun}\Op_{{}^LG}(D^\times)$ 中的理想由子代数 $\operatorname{Fun}\Op_{{}^LG}(D^\times)_{<0}$ 生成，定理 8.3.1 的同构给出代数同构
\begin{equation}
    Z'\simeq\left(\operatorname{Fun}\Op_{{}^LG}(D^\times)/(\operatorname{Fun}\Op_{{}^LG}(D^\times)_{<0})\right)_0
    =(\operatorname{Fun}\Op^{\mathrm{RS}}_{{}^LG}(D))_0 .
    \tag{9.4-3}
\end{equation}
因此交换图 (9.4-2) 表明映射 (9.4-1) 是同构。\transnote{(9.4-3) 是本节的技术核心：$Z'$ 只依赖 $\Op_{{}^LG}(D^\times)$ 的 0 次部分，而把负次部分商掉之后剩下的恰是 $\Op^{\mathrm{RS}}_{{}^LG}(D)$ 的函数代数。换言之，「中心的 0 次部分」只看见 opers 的正则奇点层，这正是 §9.5 把自同态代数等同于 $\operatorname{Fun}\Op^{\mathrm{RS}}_{{}^LG}(D)_{\varpi}$ 的出发点。落点：§9.4.1，(9.4-3)，原书 p.277。[来源：本书 p.277；定理 8.3.1。]}

回忆在定理 8.3.3 的同构中有 $b_{i,n}\mapsto-u_{i,n}$。于是得到：

\begin{proposition}[{[Frenkel 2005b]}, Prop. 12.8]
    存在同构的交换图
    \[
        \begin{CD}
            Z(\mathfrak{g}) @>>> (\operatorname{Fun}\Op^{\mathrm{RS}}_{{}^LG}(D))_0\\
            @VVV @AAA\\
            (\operatorname{Fun}\mathfrak{h}^*)^W @>>> (\operatorname{Fun}\mathfrak{h}^*)^W
        \end{CD}
    \]
    其中下水平箭头由 $f\mapsto f^-$，$f^-(\lambda)=f(-\lambda)$ 给出。
\end{proposition}

\subsection{诱导函子}

我们定义一个从 $\mathfrak{g}$-模范畴到 $\mathbf{1}$ 作用为恒等的 $\gghat_{\kappa_c}$-模范畴的函子。设 $M$ 是 $\mathfrak{g}$-模。定义诱导 $\gghat_{\kappa_c}$-模 $\operatorname{Ind}(M)$ 为
\[
    \operatorname{Ind}(M)=\operatorname{Ind}^{\gghat_{\kappa_c}}_{\mathfrak{g}[[t]]\oplus\mathbf{1}}M,
\]
其中 $\mathfrak{g}[[t]]$ 经赋值同态 $\mathfrak{g}[[t]]\to\mathfrak{g}$ 作用在 $M$ 上，而 $\mathbf{1}$ 作用为恒等。

特别地，若 $M=M_\lambda$ 是最高权 $\lambda\in\mathfrak{h}^*$ 的 Verma 模，则诱导模 $\operatorname{Ind}(M_\lambda)$ 是临界水平 $\gghat_{\kappa_c}$ 上的 Verma 模，记作 $\mathbb{M}_{\lambda,\kappa_c}$，此前已遇到过。

% [Tutor Note: p.277 的诱导用 $\mathfrak{g}[[t]]\oplus\mathbf{1}$（只含常数项），而 §9.4.3 中的 $\hat{\mathfrak{n}}_+$ 是 $\mathfrak{n}_+\oplus t\mathfrak{g}[[t]]$。两者相差“$t\mathfrak{g}[[t]]$ 作用平凡”这一条件，这是把诱导模与 Wakimoto 模作比较的基础；混淆这两个子代数会导致 $I^0$ 的作用算错。]

再设 $Z(\mathfrak{g})$ 是 $U(\mathfrak{g})$ 的中心。Harish--Chandra 同态 $Z(\mathfrak{g})\hookrightarrow\operatorname{Fun}(\mathfrak{h}^*)^W$ 给出同构 $\mathfrak{h}^*/W\to\Spec Z(\mathfrak{g})$。于是每个 $\lambda\in\mathfrak{h}^*$ 定义一个点 $\varpi(\lambda)\in\mathfrak{h}^*/W$，从而定义 $Z(\mathfrak{g})$ 的一个特征标，我们也记作 $\varpi(\lambda)$。特别地，$Z(\mathfrak{g})$ 经中心特征标 $\varpi(\lambda)$ 作用在 Verma 模 $M_{\lambda-\rho}$ 上。

\begin{proposition}
    设 $M$ 是任意 $\mathfrak{g}$-模。把 $\operatorname{Ind}(M)$ 看成完备泛包络代数 $\widetilde{U}_{\kappa_c}(\gghat)$ 的中心 $Z(\gghat)$ 上的模。则 $\operatorname{Ind}(M)$ 在概形论意义下支撑在子概形 $\Op^{\mathrm{RS}}_{{}^LG}(D)\subset\Op_{{}^LG}(D^\times)$ 上，即它被 $Z(\gghat)\simeq\operatorname{Fun}\Op_{{}^LG}(D^\times)$ 中 $\Op^{\mathrm{RS}}_{{}^LG}(D)$ 的理想湮灭。

    此外，若 $Z(\mathfrak{g})$ 经特征标 $\varpi(\lambda)$ 作用在 $M$ 上，则 $\operatorname{Ind}(M)$ 在概形论意义下支撑在子概形 $\Op^{\mathrm{RS}}_{{}^LG}(D)_{\varpi(-\lambda)}$ 上。\transnote{「概形论意义下支撑」比「作为模有某个中心特征标」强：它要求 $\operatorname{Ind}(M)$ 被 $\Op^{\mathrm{RS}}_{{}^LG}(D)_{\varpi(-\lambda)}$ 的理想（而不只是被 $\Op^{\mathrm{RS}}_{{}^LG}(D)$ 的理想）湮灭。这里出现 $-\lambda$ 而非 $\lambda$，是因为 $M$ 落在 0 次部分，而留数映射与 Harish--Chandra 同态的复合带一个 $f\mapsto f^-$、$f^-(\lambda)=f(-\lambda)$ 的扭转（命题 9.4.1）。落点：§9.4.2 命题 9.4.2，原书 p.278。[来源：本书 p.278 与命题 9.4.1。]}
\end{proposition}

\begin{proof}
    这样定义 $L_0=-t\partial_t$ 的作用：它在生成子空间 $M\subset\operatorname{Ind}(M)$ 上作用为 0，并与 $L_0$ 在 $\gghat_{\kappa_c}$ 上的作用相容。于是 $\operatorname{Ind}(M)$ 获得 $\Z$-分次，它在 $\operatorname{Ind}(M)$ 上只取非负值，0 次部分恰为子空间 $M$。这说明 $Z(\gghat)$ 中由一切负次元素张成的子代数 $Z(\gghat)_{<0}$ 在 $\operatorname{Ind}(M)$ 上作用为 0。算子 $L_0$ 也在 $\operatorname{Fun}\Op_{{}^LG}(D^\times)$ 上定义 $\Z$-分次，且同构 $Z(\gghat)\simeq\operatorname{Fun}\Op_{{}^LG}(D^\times)$ 与该分次相容。因此在此同构下 $Z(\gghat)_{<0}$ 映入一切负次元素张成的子代数 $\operatorname{Fun}\Op_{{}^LG}(D^\times)_{<0}$。

    但 $\operatorname{Fun}\Op_{{}^LG}(D^\times)_{<0}$ 在 $\operatorname{Fun}\Op_{{}^LG}(D^\times)$ 中的理想，恰是子空间 $\Op^{\mathrm{RS}}_{{}^LG}(D)\subset\Op_{{}^LG}(D^\times)$（即具有正则奇点的 opers）的理想，如 §9.4.1 所示。这给出命题的第一个断言。

    其次考虑 $Z(\gghat)$ 的 0 次子代数 $Z(\gghat)_0$ 在 $\operatorname{Ind}(M)$ 上的作用。它经由商
    \[
        Z'=Z(\gghat)_0/(Z(\gghat)\cdot Z(\gghat)_{<0})_0
    \]
    分解。由命题 9.4.1，$Z'$ 同构于 $U(\mathfrak{g})$ 的中心 $Z(\mathfrak{g})$，且它在 $M\subset\operatorname{Ind}(M)$ 上的作用经此同构分解。此外，由命题 9.4.1 有交换图
    \[
        \begin{CD}
            Z' @>>> \operatorname{Fun}\Op^{\mathrm{RS}}_{{}^LG}(D)\\
            @VVV @AAA\\
            (\operatorname{Fun}\mathfrak{h}^*)^W @>>> (\operatorname{Fun}\mathfrak{h}^*)^W
        \end{CD}
    \]
    其中左竖箭头是同构 $Z'\simeq Z(\mathfrak{g})$ 与 Harish--Chandra 同态的复合，右竖箭头由 §9.1.1 的留数映射
    \[
        \operatorname{res}:\Op^{\mathrm{RS}}_{{}^LG}(D)\to\mathfrak{h}^*/W
    \]
    诱导，而下水平箭头由 $f\mapsto f^-$，$f^-(\lambda)=f(-\lambda)$ 给出。这给出命题的第二个断言。
\end{proof}

特别地，Verma 模 $\mathbb{M}_{\mu,\kappa_c}$ 在概形论意义下支撑在子概形 $\Op^{\mathrm{RS}}_{{}^LG}(D)_{\varpi(-\mu-\rho)}$ 上。

% [Tutor Note: p.278–279 的“概形论意义下支撑”是本节统一语言：中心 $Z(\gghat)\simeq\operatorname{Fun}\Op_{{}^LG}(D^\times)$ 的作用经某个理想分解，就等价于说模“支撑”在对应子概形上。注意第二断言里的负号：$Z(\mathfrak{g})$ 的特征标是 $\varpi(\lambda)$，而支撑点是 $\varpi(-\lambda)$，来自 Miura 变换换到 ${}^LG$-侧时留数变号（见 p.279 脚注 2）。此符号是最易出错处。]

\subsection{Wakimoto 模与 $\gghat$-模的范畴}

考虑李代数
\[
    \hat{\mathfrak{n}}_+ = (\mathfrak{n}_+\otimes1)\oplus(\mathfrak{g}\otimes t\C[[t]]),
\]
它是 $\mathfrak{g}((t))$ 与 $\gghat_{\kappa_c}$ 的子代数。相应地，$G((t))$ 的子群记作 $I^0$：它是李子群 $N_+\subset G$（对应 $\mathfrak{n}_+\subset\mathfrak{g}$）在赋值映射 $G[[t]]\to G$ 下的原像。注意 $I^0=[I,I]$，其中 $I$ 是 §10.2.2 中定义的 Iwahori 子群。

对中心特征标 $\varpi(\nu)\in\Spec Z(\mathfrak{g})\simeq\mathfrak{h}^*/W$，令 $\mathcal{O}_{\varpi(\nu)}$ 为 $\mathfrak{g}$-模范畴 $\mathcal{O}$（关于固定的 Borel 子代数 $\mathfrak{b}$）中使 $Z(\mathfrak{g})$ 经特征标 $\varpi(\nu)$ 作用的模所成的块。令 $\gghat_{\kappa_c}\text{-mod}^{I^0}_{\varpi(-\nu)}$ 为满足以下条件的 $\gghat_{\kappa_c}$-模 $M$ 所成的范畴：
\begin{itemize}
    \item 李代数 $\hat{\mathfrak{n}}_+$ 在 $M$ 上的作用局部幂零；
    \item $Z(\gghat)$ 在 $M$ 上的作用经由 $\operatorname{Fun}\Op^{\mathrm{RS}}_{{}^LG}(D)_{\varpi(-\nu)}$ 分解。
\end{itemize}
在这样的模上，$\hat{\mathfrak{n}}_+$ 的作用可指数化为 $I^0$ 的作用，记号的来源即在此。\transnote{两个范畴都是「块」：$\mathcal{O}_{\varpi(\nu)}$ 是有限维范畴 $\mathcal{O}$ 中中心特征标固定为 $\varpi(\nu)$ 的块，$\gghat_{\kappa_c}\text{-mod}^{I^0}_{\varpi(-\nu)}$ 则是临界水平的对应物，其定义的两条正是块的两个要件——$\hat{\mathfrak{n}}_+$ 局部幂零（保证可积到 $I^0$）与中心作用经 $\operatorname{Fun}\Op^{\mathrm{RS}}_{{}^LG}(D)_{\varpi(-\nu)}$ 分解（保证中心特征标落在单层上）。用 $I^0$ 而不用 $I$，是因为 $\hat{\mathfrak{n}}_+$ 的李代数是 $\mathfrak{n}_+\oplus t\mathfrak{g}[[t]]$，指数化只得到 $I^0$；但指标写成 $-\nu$ 与 $I^0$ 无关，那是命题 9.4.2 中留数变号的直接后果。落点：§9.4.3，原书 p.279。[来源：本书 p.279 与命题 9.4.2。]}

于是由命题 9.4.2，诱导函子 $\operatorname{Ind}$ 给出函子
\[
    \operatorname{Ind}:\mathcal{O}_{\varpi(\nu)}\to\gghat_{\kappa_c}\text{-mod}^{I^0}_{\varpi(-\nu)}.
\]
特别地，对任意 $\lambda\in\mathfrak{h}^*$，Verma 模 $\mathbb{M}_{w(\lambda+\rho)-\rho,\kappa_c}$（$w\in W$）是范畴 $\gghat_{\kappa_c}\text{-mod}^{I^0}_{\varpi(-\lambda-\rho)}$ 的对象。此外容易看出，$\gghat_{\kappa_c}\text{-mod}^{I^0}_{\varpi(-\lambda-\rho)}$ 的任一对象都有滤过，其分次商是这些 Verma 模的商。

现在要说明临界水平的某些 Wakimoto 模也是范畴 $\gghat_{\kappa_c}\text{-mod}^{I^0}_{\varpi(-\lambda-\rho)}$ 的对象。这些模的一般构造以及中心 $Z(\gghat)$ 在它们上面的相应作用，已在定理 8.3.3 中描述。

回忆 §9.3.1 中的子空间 $\operatorname{Conn}(\Omega^\rho)^{-\lambda}_D\subset\operatorname{Conn}(\Omega^\rho)_{D^\times}$。在同构 $\operatorname{Conn}(\Omega^\rho)_{D^\times}\simeq\operatorname{Conn}(\Omega^{-\rho})_{D^\times}$ 下，它变成由形如
\begin{equation}
    \partial_t+\frac{\lambda}{t}+\mathbf{u}(t),\qquad \mathbf{u}(t)\in{}^L\mathfrak{h}[[t]]
    \tag{9.4-4}
\end{equation}
的联络所成的子空间 $\operatorname{Conn}(\Omega^{-\rho})^\lambda_D$。\transnote{这里的 $\lambda/t$ 项就是留数：联络 $\partial_t+\lambda/t+\mathbf{u}(t)$ 在 $t=0$ 处的留数是 $\varpi(\lambda)$，故「连通部分张成的仿射空间 $\operatorname{Conn}(\Omega^{-\rho})^\lambda_D$」由留数 $\lambda$ 决定。同构 $\operatorname{Conn}(\Omega^{\rho})_{D^\times}\simeq\operatorname{Conn}(\Omega^{-\rho})_{D^\times}$ 换丛时 $u\mapsto-u$、$\lambda\mapsto-\lambda$，故同一子空间既可写成 $\operatorname{Conn}(\Omega^\rho)^{-\lambda}_D$ 又可写成 $\operatorname{Conn}(\Omega^{-\rho})^\lambda_D$；Miura 变换把留数 $\lambda$ 送到 $-\lambda-\rho$（换丛变号加上 $\rho$-平移），于是像所在层的指标记作 $\varpi(-\lambda-\rho)$。三处符号不同源，读下文须逐一核对。落点：§9.4.3，(9.4-4)，原书 p.279–280。[来源：本书 p.279–280 与 §9.3.1。]}

% [Tutor Note: p.279 式 (9.4-4) 是本节全部构造的坐标模型：参数 $\lambda$ 与 $\lambda/t$ 项一起决定正则奇点的留数。注意 $D$ 与 $D^\times$ 之别：$\operatorname{Conn}(\Omega^{-\rho})^\lambda_D$ 由在闭圆盘上无极点的联络组成，是仿射空间，其函数环是多项式环（见 (9.6-10)），而 $\operatorname{Conn}(\Omega^{-\rho})_{D^\times}$ 上的函数环要大得多。]

注意 $\operatorname{Fun}\operatorname{Conn}(\Omega^{-\rho})^\lambda_D$ 是代数 $\operatorname{Fun}\operatorname{Conn}(\Omega^{-\rho})_{D^\times}$ 上的光滑模。事实上，$\operatorname{Fun}\operatorname{Conn}(\Omega^{-\rho})^\lambda_D$ 自然地与交换顶点代数 $\pi_0$ 上的模 $\pi^0_\lambda$ 等同（见 §6.2.1）。因此我们在
\begin{equation}
    W_\lambda\overset{\mathrm{def}}{=}M_{\mathfrak{g}}\otimes\operatorname{Fun}\operatorname{Conn}(\Omega^{-\rho})^\lambda_D
    \tag{9.4-5}
\end{equation}
上得到 $\gghat_{\kappa_c}$-模结构。

我们断言 $W_\lambda$ 是范畴 $\gghat_{\kappa_c}\text{-mod}^{I^0}_{\varpi(-\lambda-\rho)}$ 的对象。

事实上，用定理 6.1.6 中给出的 $\gghat_{\kappa_c}$ 作用的显式公式可以验证：李代数 $\hat{\mathfrak{n}}_+$ 在 $W_\lambda$ 上局部幂零地作用。另一方面，由 §9.3.1，Miura 变换在 $\operatorname{Conn}(\Omega^{-\rho})^\lambda_D=\operatorname{Conn}(\Omega^\rho)^{-\lambda}_D$ 上的限制取值于
$\Op^{\mathrm{RS}}_{{}^LG}(D)_{\varpi(-\lambda-\rho)}\subset\Op_{{}^LG}(D^\times)$。因此由定理 8.3.3，$Z(\gghat)$ 在 $W_\lambda$ 上的作用经由商 $\operatorname{Fun}\Op^{\mathrm{RS}}_{{}^LG}(D)_{\varpi(-\lambda-\rho)}$ 分解。故 $W_\lambda$ 是范畴 $\gghat_{\kappa_c}\text{-mod}^{I^0}_{\varpi(-\lambda-\rho)}$ 的对象。

这说明 $W_\lambda$ 的任何商也是 $\gghat_{\kappa_c}\text{-mod}^{I^0}_{\varpi(-\lambda-\rho)}$ 的对象。代数 $\operatorname{Fun}\operatorname{Conn}(\Omega^{-\rho})^\lambda_D$ 通过自同态作用在 $W_\lambda$ 上，这些自同态与 $\gghat_{\kappa_c}$ 的作用交换。给定 $\overline{\nabla}\in\operatorname{Conn}(\Omega^{-\rho})^\lambda_D$，令 $W_{\overline{\nabla}}$ 为 $W_\lambda$ 模掉 $\operatorname{Fun}\operatorname{Conn}(\Omega^{-\rho})^\lambda_D$ 中对应于 $\overline{\nabla}$ 的极大理想所得的商。这正是由 $\operatorname{Fun}\operatorname{Conn}(\Omega^{-\rho})_{D^\times}$ 上对应于 $\overline{\nabla}\in\operatorname{Conn}(\Omega^{-\rho})^\lambda_D\subset\operatorname{Conn}(\Omega^{-\rho})_{D^\times}$ 的一维模 $\C_{\overline{\nabla}}$ 得到的 Wakimoto 模 $M_{\mathfrak{g}}\otimes\C_{\overline{\nabla}}$。由定理 8.3.3，中心 $Z(\gghat)$ 经中心特征标 $\mu(\overline{\nabla})$ 作用在 $W_{\overline{\nabla}}$ 上，其中 $\mu$ 是 Miura 变换。

对 $\chi\in\Op^{\mathrm{RS}}_{{}^LG}(D)$，令 $\gghat_{\kappa_c}\text{-mod}^{I^0}_\chi$ 为满足以下条件的 $\gghat_{\kappa_c}$-模 $M$ 所成的范畴：
\begin{itemize}
    \item $\hat{\mathfrak{n}}_+$ 在 $M$ 上的作用局部幂零；
    \item 中心 $Z(\gghat)$ 经对应于 $\chi$ 的特征标 $Z(\gghat)\to\C$ 作用在 $M$ 上。
\end{itemize}
设 $\chi$ 的留数为 $\varpi(-\lambda-\rho)\in\mathfrak{h}^*/W$（$\lambda\in\mathfrak{h}^*$）。则 Verma 模 $\mathbb{M}_{w(\lambda+\rho)-\rho,\kappa_c}$（$w\in W$）模掉对应于 $\chi$ 的中心特征标所得的商 $\mathbb{M}_{w(\lambda+\rho)-\rho,\kappa_c}(\chi)$ 是该范畴的对象。此外容易看出，$\gghat_{\kappa_c}\text{-mod}^{I^0}_\chi$ 的任一对象都有滤过，其分次商是 $\mathbb{M}_{w(\lambda+\rho)-\rho,\kappa_c}(\chi)$ 的商。

Wakimoto 模 $W_{\overline{\nabla}}$ 是范畴 $\gghat_{\kappa_c}\text{-mod}^{I^0}_\chi$ 的对象，当且仅当 $\overline{\nabla}\in\mu^{-1}(\chi)$。由同构 (9.3-6)（其中把 $G$ 换成 ${}^LG$、把 $\Omega^{\check\rho}$ 换成 $\Omega^{-\rho}$，从而改变联络留数的符号），这意味着
\begin{equation}
    \overline{\nabla}\in\operatorname{Conn}(\Omega^{-\rho})^{w(\lambda+\rho)-\rho}_D,
    \qquad w\in W.
    \tag{9.4-6}
\end{equation}

% [Tutor Note: p.280 的 $W_{\overline{\nabla}}$ 由 $\operatorname{Conn}$ 空间上的极大理想模掉得到，是“以连接为参数”的 Wakimoto 模族；它与 §9.5 中带 $w$ 上标的族 $W^w_\lambda$ 是两个不同的构造（前者对应 $w=1$）。\ \S 9.4 与 §9.5 的 $\operatorname{Conn}$ 上标（$\lambda$ 为参数）与 $w(\lambda+\rho)-\rho$ 之间的换算，全靠 (9.4-6) 这一步。]

Cartan 子代数 $\mathfrak{h}\otimes1$ 对角地作用在 $W_{\overline{\nabla}}$ 上；但它的作用可指数化为相应群 $H$ 的作用（从而 $W_{\overline{\nabla}}$ 是 $I$-等变的），当且仅当 $\lambda$ 属于整优权集合 $P^+$。

于是设 $\lambda\in P^+$。此时
\[
    \chi\in\Op^{\mathrm{RS}}_{{}^LG}(D)_{\varpi(-\lambda-\rho)}=\Op^{\mathrm{nilp},\lambda}_{{}^LG}
\]
（见 §9.2.1）。相应的范畴 $\gghat_{\kappa_c}\text{-mod}^{I^0}_\chi$ 将在 §10.4 讨论。它们将作为圈群局部 Langlands 对应赋予 $D^\times$ 上驯顺分歧局部系统（tamely ramified local systems）的范畴出现。下面这个事实在此语境中对我们将非常重要（见 §10.4.5）。

由定理 9.3.7，$\mu^{-1}(\chi)$ 的点与 $\operatorname{Res}_{\check\lambda}(\chi)$ 的 Springer 纤维 $\operatorname{Sp}_{\operatorname{Res}_{\check\lambda}(\chi)}$ 的点一一对应。因此，Wakimoto 模 $W_{\overline{\nabla}}$（$\overline{\nabla}\in\mu^{-1}(\chi)$）给出范畴 $\gghat_{\kappa_c}\text{-mod}^{I^0}_\chi$ 中以 $\operatorname{Sp}_{\operatorname{Res}_{\check\lambda}(\chi)}$ 的点为参数的一族对象。

特别地，若 $\operatorname{Res}_{\check\lambda}(\chi)=0$，从而 $\chi\in\Op^{\check\lambda}_{{}^LG}$，则我们得到范畴 $\gghat_{\kappa_c}\text{-mod}^{I^0}_\chi$ 中以旗簇 ${}^LG/{}^LB$（即 Springer 纤维 $\operatorname{Sp}_0$）的点为参数的一族 Wakimoto 模。

由此可见，当
\[
    \chi\in\Op^{\mathrm{RS}}_{{}^LG}(D)_{\varpi(-\lambda-\rho)},
\]
且 $\lambda\in P^+$ 时，范畴 $\gghat_{\kappa_c}\text{-mod}^{I^0}_\chi$ 相当复杂。这是范畴 $\gghat_{\kappa_c}\text{-mod}^{I^0}_\chi$ 的一个“极端”例子。另一个“极端”是范畴 $\gghat_{\kappa_c}\text{-mod}^{I^0}_\chi$，其中 $\chi$ 同上，而 $\lambda$ 是 $\mathfrak{h}^*$ 中满足
\[
    \langle\alpha,w(\check\lambda+\rho)-\rho\rangle\notin\Z_+,
    \qquad \alpha\in\Delta_+,\ w\in W
\]
的一般元素。此时该范畴的结构相当简单。应用命题 9.3.1（其中再次把 $G$ 换成 ${}^LG$、把 $\Omega^{\check\rho}$ 换成 $\Omega^{-\rho}$，从而改变留数符号），得到
\[
    \mu^{-1}(\chi)\cap\operatorname{Conn}(\Omega^{-\rho})^{w(\lambda+\rho)-\rho}_D
\]
由单点组成，记作 $\overline{\nabla}_w$。由 (9.4-6)，这意味着此时 $\mu^{-1}(\chi)$ 的点数与 外尔群 $W$ 的元素个数相同，在每个子空间 $\operatorname{Conn}(\Omega^{-\rho})^{w(\lambda+\rho)-\rho}_D$ 中各有一个（因 $\lambda$ 一般，它自动正则）。因此此时范畴 $\gghat_{\kappa_c}\text{-mod}^{I^0}_\chi$ 有 $|W|$ 个两两不等价的 Wakimoto 模 $W_{\overline{\nabla}_w}$（$w\in W$）。

对每个 $w\in W$ 有非零同态 $\mathbb{M}_{w(\lambda+\rho)-\rho,\kappa_c}\to W_{\overline{\nabla}_w}$。在这些模是 $\Z$-分次的特殊情形（即 $\chi$ 在向量场 $L_0=-t\partial_t$ 下不变），我们在 §6.4 已证明它们全都不可约，且该同态是同构。因此对一般的 $\chi$ 同样成立（我们预期对留数为 $\varpi(-\lambda-\rho)$（$\lambda$ 如上）的一切 $\chi$ 都成立）。于是此时范畴 $\gghat_{\kappa_c}\text{-mod}^{I^0}_\chi$ 有 $|W|$ 个互不同构的不可约对象 $W_{\overline{\nabla}_w}$（$w\in W$）。分析 Cartan 子代数 $\mathfrak{h}\otimes1$ 在它们上面的作用可知，它们之间不可能有非平凡扩张。因此范畴 $\gghat_{\kappa_c}\text{-mod}^{I^0}_\chi$ 的任一对象都是若干个不可约 Wakimoto 模 $W_{\overline{\nabla}_w}$（$w\in W$）的直和。

对一般的 $\lambda\in\mathfrak{h}^*$，范畴 $\gghat_{\kappa_c}\text{-mod}^{I^0}_\chi$（$\chi\in\Op^{\mathrm{RS}}_{{}^LG}(D)_{\varpi(-\lambda-\rho)}$）的结构介于上述两种极端情形之间。

% [Tutor Note: p.281 的两个极端是理解 §10.4 的关键：$\lambda\in P^+$ 时中心特征标落在 $\Op^{\mathrm{nilp},\lambda}$ 上，Wakimoto 模的参数空间是整个 Springer 纤维（可以很奇异）；$\lambda$ 一般时 Springer 纤维退化为 $|W|$ 个单点，范畴半单。注意“驯顺分歧局部系统”这一说法在此处首次出现，它是 §10.4 中局部 Langlands 对应的靶。]

\section{Verma 模的自同态}

\begin{keypoints}
    \item \textbf{核心思想：}临界诱导模的自同态代数等于 ${}^LG$-oper 空间某子簇上的函数代数；Verma 模 $\mathbb{M}_\nu$ 对应 $\Op^{\mathrm{RS}}_{{}^LG}(D)_{\varpi(-\nu-\rho)}$，Weyl 模对应 $\Op^{\check\lambda}_{{}^LG}$。
    \item \textbf{数学工具：}$\hat{\mathfrak{n}}_+$-不变向量与最高权向量的自由度；用 $w^{-1}$-扭把 Verma 模与 Wakimoto 模 $W^w_{w(\lambda+\rho)-\rho}$ 等同（命题 9.5.1）；Kac–Kazhdan 定理控制子商最高权。
    \item \textbf{关键结论：}定理 9.5.3：$Z(\gghat)\twoheadrightarrow\operatorname{End}\mathbb{M}_\nu$，且 $\operatorname{End}\mathbb{M}_\nu\simeq\operatorname{Fun}\Op^{\mathrm{RS}}_{{}^LG}(D)_{\varpi(-\nu-\rho)}$；$\mathbb{M}_\nu$ 是该代数上的自由模。
    \item \textbf{记号提醒：}$\mathbb{M}_\nu$ 是临界水平 $\gghat_{\kappa_c}$ 上的 Verma 模（书用 blackboard 体）；$M^w_{\mathfrak{g}}$ 是 Weyl 代数 $\mathcal{A}^{\mathfrak{g}}$ 的模，而 $W^w_\lambda$ 是相应的 Wakimoto 模。
\end{keypoints}

本节与下一节把某些临界水平诱导 $\gghat$-模的自同态代数，与本章前面构造的 $\Op_{{}^LG}(D^\times)$ 的某些子簇上的函数代数等同起来。

先看真空模 $V_{\kappa_c}(\mathfrak{g})$ 的情形。由公式 (3.3-3) 有
\[
    \operatorname{End}_{\gghat_{\kappa_c}}V_{\kappa_c}(\mathfrak{g})\simeq\mathfrak{z}(\gghat).
\]
因此由定理 4.3.2 有同构
\begin{equation}
    \operatorname{End}_{\gghat_{\kappa_c}}V_{\kappa_c}(\mathfrak{g})\simeq\operatorname{Fun}\Op_{{}^LG}(D).
    \tag{9.5-1}
\end{equation}
此外，同态 $Z(\gghat)\to\operatorname{End}_{\gghat_{\kappa_c}}V_{\kappa_c}(\mathfrak{g})$ 是满的，对应同态 $\operatorname{Fun}\Op_{{}^LG}(D^\times)\to\operatorname{Fun}\Op_{{}^LG}(D)$。于是 $V_{\kappa_c}(\mathfrak{g})$ 对应子簇
$\Op_{{}^LG}(D)\subset\Op_{{}^LG}(D^\times)$。

这里把该结果推广到 Verma 模的情形，它们对应子簇 $\Op^{\mathrm{RS}}_{{}^LG}(D)_\varpi$；下一节推广到 Weyl 模的情形，它们对应 $\Op^{\check\lambda}_{{}^LG}$。

\subsection{Wakimoto 模的族}

先考虑 Verma 模。回忆水平 $\kappa$、最高权 $\lambda\in\mathfrak{h}^*$ 的 Verma 模定义为诱导模 (6.3-9)。由于此后只考虑临界水平 $\kappa=\kappa_c$，我们从记号中去掉下标 $\kappa_c$，简写为 $\mathbb{M}_\lambda$。

注意 $\mathbb{M}_\lambda=\operatorname{Ind}(M_\lambda)$，其中 $\operatorname{Ind}$ 是 §9.4.2 引入的诱导函子。

我们希望描述 $\mathbb{M}_\lambda$ 的自同态代数，对一切最高权 $\lambda\in\mathfrak{h}^*$。为此，把每个 Verma 模与某个自同态代数容易计算的 Wakimoto 模等同起来。命题 6.3.3 给出了这种等同的一个例子。该命题提示我们需用比前面考虑的模 $W_\lambda=M_{\mathfrak{g}}\otimes\pi_\lambda$ 更一般的 Wakimoto 模（例如见 §9.4.3）。确切地说，我们将使用 §6.3 的半无限抛物诱导、关于 $\mathfrak{g}$ 的形如 $w\mathfrak{b}_-w^{-1}$ 的 Borel 子代数所得的 Wakimoto 模，其中 $\mathfrak{b}_-$ 是我们此前一直使用的标准 Borel 子代数，$w$ 是 $\mathfrak{g}$ 的外尔群元素。

本节对 $\mathfrak{g}$ 的外尔群每个元素 $w$ 定义一族 Wakimoto 模。它把族 $W_\lambda$（对应 $w=1$，见 §9.4.3）的构造按 §6.3 的一般构造加以推广。

回忆 §5.3.3 引入的外尔代数 $\mathcal{A}^{\mathfrak{g}}$。定义 $\mathcal{A}^{\mathfrak{g}}$-模 $M^w_{\mathfrak{g}}$ 为由真空向量 $|0\rangle_w$ 生成、且满足
\[
    \begin{aligned}
        &a_{\alpha,n}|0\rangle_w=a^*_{\alpha,n}|0\rangle_w=0, && n>0;\\
        &a_{\alpha,0}|0\rangle_w=0, && \alpha\in\Delta_+\cap w^{-1}(\Delta_+);\\
        &a^*_{\alpha,0}|0\rangle_w=0, && \alpha\in\Delta_+\cap w^{-1}(\Delta_-)
    \end{aligned}
\]
的模。由定理 8.3.3，对每个 $\lambda\in\mathfrak{h}^*$，张量积
$M^w_{\mathfrak{g}}\otimes\operatorname{Fun}\operatorname{Conn}(\Omega^{-\rho})^\lambda_D$ 是 $\gghat_{\kappa_c}$-模。记相应的同态为 $\omega$
\[
    \gghat_{\kappa_c}\to\operatorname{End}\left(M^w_{\mathfrak{g}}\otimes\operatorname{Fun}\operatorname{Conn}(\Omega^{-\rho})^\lambda_D\right).
\]
定义一个新的 $\gghat_{\kappa_c}$-模，记作 $W^w_\lambda$，为 $M^w_{\mathfrak{g}}\otimes\operatorname{Fun}\operatorname{Conn}(\Omega^{-\rho})^\lambda_D$ 的 $w^{-1}$-扭。确切地说，用公式 $x\mapsto\omega(\operatorname{Ad}\widetilde{w}^{-1}\cdot x)$ 定义 $\gghat_{\kappa_c}$ 的新作用，其中 $\widetilde{w}$ 是 $w$ 到群 $G$ 的 Tits 提升。\transnote{扭的方向值得单独交代：$M^w_{\mathfrak{g}}$ 上的原始作用是 $[h,a_{\alpha,n}]=\langle\alpha,h\rangle a_{\alpha,n}$，用 $w$ 扭之后变成 $[h,a_{\alpha,n}]=\langle w(\alpha),h\rangle a_{\alpha,n}$（即 (9.5-3)），于是 $|0\rangle_w\otimes1$ 的权从 $\lambda$ 变成 $w(\lambda+\rho)-\rho$，这正是把 Verma 模的最高权搬到 $w$ 的轨道上所需的一步。$w=1$ 时扭是恒等，$W^1_\lambda=W_\lambda$ 可作核对；$\widetilde{w}$ 换成 $z\widetilde{w}$（$z$ 在中心）不改变扭后的作用。落点：§9.5.1，原书 p.283。[来源：本书 p.283 与式 (9.5-3)。]}

特别地，$W^1_\lambda$ 就是 §9.4.3 引入的模 $W_\lambda$。

回忆：给定 $\gghat_{\kappa_c}$-模，若算子 $L_0=-t\partial_t$ 与李子代数 $\mathfrak{h}\otimes1\subset\gghat_{\kappa_c}$ 在其上对角作用且公共本征空间有限维，则定义 $M$ 的特征标为形式级数
\begin{equation}
    \operatorname{ch}M=\sum_{\hat\lambda\in(\mathfrak{h}\oplus\C L_0)^*}\dim M(\hat\lambda)\,e^{\hat\lambda},
    \tag{9.5-2}
\end{equation}
其中 $M(\hat\lambda)$ 是 $L_0$ 与 $\mathfrak{h}\otimes1$ 对应于权 $\hat\lambda:(\mathfrak{h}\oplus\C L_0)\to\C$ 的本征空间。

% [Tutor Note: p.283 的 $w^{-1}$-扭定义（$x\mapsto\omega(\operatorname{Ad}\widetilde{w}^{-1}\cdot x)$）极易写反。核对办法：$w=1$ 时须退回 $W^1_\lambda=W_\lambda$；再取 $w=w_0$、$\lambda=-2\rho$，须与命题 6.3.3 的 $W^+_{0,\kappa_c}$ 一致。$\widetilde{w}$ 的选取只差中心，不影响扭后的表示。]

下面计算模 $W^w_\lambda$ 的特征标。$M^w_{\mathfrak{g}}\otimes\operatorname{Fun}\operatorname{Conn}(\Omega^{-\rho})^\lambda_D$ 上 $L_0$ 的作用是 $M^w_{\mathfrak{g}}$ 与 $\operatorname{Fun}\operatorname{Conn}(\Omega^{-\rho})^\lambda_D$ 上相应算子作用之和。前者由对易关系
\[
    [L_0,a_{\alpha,n}]=-na_{\alpha,n},\qquad
    [L_0,a^*_{\alpha,n}]=-na^*_{\alpha,n}
\]
以及条件 $L_0|0\rangle_w=0$ 确定。后者由向量场 $-t\partial_t$ 在空间 $\operatorname{Conn}(\Omega^{-\rho})^\lambda_D$ 上的作用确定。若把 $\operatorname{Conn}(\Omega^{-\rho})^\lambda_D$ 的元素写成 (9.4-4) 的形式，并把 $u_i(t)=\langle\alpha_i,\mathbf{u}(t)\rangle$ 展成级数
\[
    u_i(t)=\sum_{n<0}u_{i,n}t^{-n-1},
\]
就得到等同
\[
    \operatorname{Fun}\operatorname{Conn}(\Omega^{-\rho})^\lambda_D\simeq\C[u_{i,n}]_{i=1,\dots,\ell;n<0}.
\]
此时 $L_0$ 就是该多项式代数上由公式
\[
    L_0\cdot u_{i,n}=-nu_{i,n}
\]
完全确定的导子。

接着考虑 $\mathfrak{h}=\mathfrak{h}\otimes1\subset\gghat_{\kappa_c}$ 在 $M^w_{\mathfrak{g}}$ 上的作用。由定理 6.1.6 中给出的原始公式可知，对任意 $h\in\mathfrak{h}$ 有
\[
    [h,a_{\alpha,n}]=\langle\alpha,h\rangle a_{\alpha,n},\qquad
    [h,a^*_{\alpha,n}]=-\langle\alpha,h\rangle a^*_{\alpha,n}.
\]
把 $\gghat_{\kappa_c}$（从而 $\mathfrak{h}$）的作用用 $w$ 扭之后，得到新的关系
\begin{equation}
    [h,a_{\alpha,n}]=\langle w(\alpha),h\rangle a_{\alpha,n},\qquad
    [h,a^*_{\alpha,n}]=-\langle w(\alpha),h\rangle a^*_{\alpha,n}.
    \tag{9.5-3}
\end{equation}
注意对任意 $\mu\in\mathfrak{h}^*$ 有 $\langle\mu,w^{-1}(h)\rangle=\langle w(\mu),h\rangle$。

为计算 $\mathfrak{h}$ 在 $|0\rangle_w\otimes1\in M^w_{\mathfrak{g}}\otimes1$ 上的作用，注意按定理 6.1.6 中 $h_{i,0}$ 的作用公式，
\begin{align*}
    h_{i,0}\cdot|0\rangle_w\otimes1
    &=\left(\langle w(\lambda),h_i\rangle-\sum_{\beta\in\Delta_+}\langle w(\beta),h_i\rangle a^*_{\beta,0}a_{\beta,0}\right)|0\rangle_w\otimes1\\
    &=\left(\langle w(\lambda),h_i\rangle-\sum_{\beta\in\Delta_+\cap w^{-1}(\Delta_-)}\langle w(\beta),h_i\rangle\right)|0\rangle_w\otimes1\\
    &=\langle w(\lambda)-w(w^{-1}(\rho)-\rho),h_i\rangle|0\rangle_w\otimes1\\
    &=\langle w(\lambda+\rho)-\rho,h_i\rangle|0\rangle_w\otimes1.
\end{align*}
因此向量 $|0\rangle_w\otimes1$ 的 $\mathfrak{h}$-权等于 $w(\lambda+\rho)-\rho$。

% [Tutor Note: p.284 这四行推导是本节符号的心脏：第二行只保留 $\beta\in\Delta_+\cap w^{-1}(\Delta_-)$，是因为 $a^*_{\beta,0}a_{\beta,0}|0\rangle_w$ 只有在 $\beta\in w^{-1}(\Delta_+)$ 时才非零（见 $M^w_{\mathfrak{g}}$ 的定义）；第三行用 $\sum_{\beta\in\Delta_+\cap w^{-1}(\Delta_-)}\langle w(\beta),h_i\rangle=\langle w(w^{-1}(\rho)-\rho),h_i\rangle$，其中 $\rho$ 是正根半和。若把 $\rho$ 的半和定义改成“正余根半和”，这一步整体差号。]

综合以上，得到模 $W^w_\lambda$ 的特征标公式
\begin{equation}
    \operatorname{ch}W^w_\lambda
    =e^{w(\lambda+\rho)-\rho}\prod_{\alpha\in\hat\Delta_+}(1-e^{-\alpha})^{-\operatorname{mult}(\alpha)},
    \tag{9.5-4}
\end{equation}
其中 $\hat\Delta_+$ 是 $\gghat$ 的正根集。这里实根对应的项来自 $\mathcal{A}^{\mathfrak{g}}$ 的生成元 $a_{\alpha,n}$ 与 $a^*_{\alpha,n}$，虚根对应的项来自 $\operatorname{Fun}\operatorname{Conn}(\Omega^{-\rho})^\lambda_D$ 的生成元 $u_{i,n}$。于是 $W^w_\lambda$ 的特征标等于最高权为 $w(\lambda+\rho)-\rho$ 的 Verma 模 $\mathbb{M}_{w(\lambda+\rho)-\rho}$ 的特征标，后者的特征标已在命题 6.3.3 的证明过程中算出。

\subsection{Verma 模与 Wakimoto 模}

现在注意，$\mathfrak{g}$ 的任意权 $\nu$ 都可以（可能以不同方式）写成 $\nu=w(\lambda+\rho)-\rho$，其中 $\lambda$ 满足：对任意 $\alpha\in\Delta_+$ 有 $\langle\lambda,\check\alpha\rangle\notin\Z_+$。注意若 $\nu+\rho$ 在正则整优权的外尔群轨道中，则 $w$ 由 $\nu$ 唯一确定。\transnote{条件 $\langle\lambda,\check\alpha\rangle\notin\Z_+$（对一切 $\alpha\in\Delta_+$）是「留数固定」的技术形式：它保证 $\nu+\rho$ 落在正则整优权的轨道上，从而分解 $\nu=w(\lambda+\rho)-\rho$ 中的 $w$ 由 $\nu$ 唯一确定。有了唯一性，Verma 模 $\mathbb{M}_\nu$ 才能被安置到单层 $\Op^{\mathrm{RS}}_{{}^LG}(D)_{\varpi(-\nu-\rho)}$ 上，命题 9.5.1 与定理 9.5.3 中出现的才都是同一个纤维而不是一族。若去掉该条件，同一 $\nu$ 会对应多个 $w$，Wakimoto 模也须按 $W$-轨道整族出现。落点：§9.5.2，原书 p.285–286。[来源：本书 p.285–286；[Kac and Kazhdan 1979]（本书 p.286 引用）。]}

下面的断言建立 Verma 模 $\mathbb{M}_{w(\lambda+\rho)-\rho}$ 与某类 Wakimoto 模之间的同构。

\begin{proposition}
    设权 $\lambda$ 满足：对一切 $\alpha\in\Delta_+$ 有 $\langle\lambda,\check\alpha\rangle\notin\Z_+$；设 $w$ 是 $\mathfrak{g}$ 的外尔群的任意元素。则 Verma 模 $\mathbb{M}_{w(\lambda+\rho)-\rho}$ 同构于 Wakimoto 模 $W^w_{w(\lambda+\rho)-\rho}$。
\end{proposition}

\begin{proof}
    当 $w=w_0$、$\lambda=-2\rho$ 时，这已在命题 6.3.3 中证明（注意 $W^+_{0,\kappa_c}=W^{w_0}_{-2\rho}$）。一般情形用类似的论证。

    首先注意，特征标公式表明向量 $|0\rangle_w\otimes1\in W^w_{w(\lambda+\rho)-\rho}$ 是最高权为 $w(\lambda+\rho)-\rho$ 的最高权向量。因此有典范同态
    \[
        \mathbb{M}_{w(\lambda+\rho)-\rho}\to W^w_{w(\lambda+\rho)-\rho},
    \]
    把 $\mathbb{M}_{w(\lambda+\rho)-\rho}$ 的最高权向量送到 $|0\rangle_w\otimes1\in W^w_{w(\lambda+\rho)-\rho}$。由于 $\mathbb{M}_{w(\lambda+\rho)-\rho}$ 与 $W^w_{w(\lambda+\rho)-\rho}$ 的特征标相等，只要证明此同态满，等价地证明 $W^w_{w(\lambda+\rho)-\rho}$ 由它的最高权向量 $|0\rangle_w\otimes1$ 生成，命题即得证。

    设 $W^w_{w(\lambda+\rho)-\rho}$ 不由最高权向量生成。则存在 $W^w_{w(\lambda+\rho)-\rho}$ 上的齐次线性泛函，其权小于最高权 $(\lambda,0)\in\mathfrak{h}^*\oplus(\C L_0)^*$，且在李代数
    \[
        \hat{\mathfrak{n}}_-=\mathfrak{n}_-\otimes\C[t^{-1}]\oplus\mathfrak{b}_+\otimes\C[t^{-1}]
    \]
    下不变，从而在其子代数
    \[
        L_-\mathfrak{n}^w_+=(w\mathfrak{n}_+w^{-1})\cap\mathfrak{n}_-\otimes\C[t^{-1}]
        \oplus(w\mathfrak{n}_+w^{-1})\cap\mathfrak{n}_+\otimes t^{-1}\C[t^{-1}]
    \]
    下不变。因此该泛函经由 $W^w_{w(\lambda+\rho)-\rho}$ 的 $L_-\mathfrak{n}^w_+$-余不变空间分解。但由 $\gghat_{\kappa_c}$ 在 $W^w_{w(\lambda+\rho)-\rho}$ 上作用的公式可知，李代数 $L_-\mathfrak{n}^w_+$ 的作用是自由的，余不变空间同构于
    \[
        \C[a^*_{\alpha,n}]_{\alpha\in\Delta_+,n\le0}
        \otimes\operatorname{Fun}\operatorname{Conn}(\Omega^{-\rho})^{w_0(\lambda+\rho)-\rho}_D .
    \]
    我们的泛函必须在该子空间的某个权向量上取非零值，且该权不等于最高权 $w(\lambda+\rho)-\rho$。用公式 (9.5-3) 可知 $a^*_{\alpha,n}$ 的权为 $n\delta-w(\alpha)$。因此该子空间中向量的权形如
    \begin{equation}
        (w(\lambda+\rho)-\rho)+N\delta-\sum_j\gamma_j,
        \tag{9.5-5}
    \end{equation}
    其中 $N\le0$，每个 $\gamma_j$ 是 $w(\Delta_+)$ 中的根。若这样的泛函存在，则 $W^w_{w(\lambda+\rho)-\rho}$ 有具此最高权的不可约子商。由于不可约最高权模的特征标线性无关，而 $W^w_{w(\lambda+\rho)-\rho}$ 与 $\mathbb{M}_{w(\lambda+\rho)-\rho}$ 的特征标相等，故 $\mathbb{M}_{w(\lambda+\rho)-\rho}$ 也必有最高权形如 (9.5-5) 的不可约子商。

    回忆 Kac–Kazhdan 定理 [Kac and Kazhdan 1979]，它描述 Verma 模不可约子商的最高权集合（见命题 6.3.3）。权 $\hat\mu=(\mu,n)$ 出现在 $\mathbb{M}_{\hat\nu}$（$\hat\nu=(\nu,0)$）的分解中，当且仅当 $n\le0$，且或者 $\mu=0$，或者存在权的有限序列 $\mu_0,\dots,\mu_m\in\mathfrak{h}^*$ 使得 $\mu_0=\mu$、$\mu_m=\nu$、$\mu_{i+1}=\mu_i\pm m_i\beta_i$，其中 $\beta_i$ 是某些正根、$m_i$ 是正整数，且满足
    \begin{equation}
        \langle\mu_i+\rho,\check\beta_i\rangle=m_i.
        \tag{9.5-6}
    \end{equation}
    把权固定为 $\hat\nu=(\nu,0)$ 的这类权 $\hat\mu=(\mu,n)$ 的集合记作 $S(\nu)$。由上述描述可知，若 $(\mu,n)\in S(\nu)$，则 $n\le0$，$\mu$ 属于 $\nu$ 在 $\rho$-平移作用下的 $W$-轨道，而且对任意 $\mu\in S(\nu)$，差 $\mu-\nu$ 是根的整系数线性组合。

    设权 $\lambda$ 满足命题中的条件：对任意 $\alpha\in\Delta_+$ 有 $\langle\lambda,\check\alpha\rangle\notin\Z_+$。这意味着若 $\mu=\gamma(\lambda+\rho)-\rho$ 且 $\mu-\lambda=\sum_{i\in I}n_i\alpha_i$（$n_i\in\Z$），则必然 $n_i\ge0$。因此对任意权 $(\mu,0)\in S(\lambda)$ 有 $\mu=\lambda+\sum_j\beta_j$，其中 $\beta_j\in\Delta_+$。这说明若 $\hat\mu'=(\mu',n)\in S(w(\lambda+\rho)-\rho)$，则
    \[
        \hat\mu'=(w(\lambda+\rho)-\rho)-N\delta+\sum_j\beta_j,
        \qquad \beta_j\in w(\Delta_+),\ N\le0.
    \]
    这样的权不可能形如 (9.5-5)，除非它就是 $w(\lambda+\rho)-\rho$。因此 $W^w_{w(\lambda+\rho)-\rho}$ 由其最高权向量生成，从而 $W^w_{w(\lambda+\rho)-\rho}$ 同构于 $\mathbb{M}_{w(\lambda+\rho)-\rho}$。\transnote{本证明中三处记号均照原书保留：最高权应作 $(w(\lambda+\rho)-\rho,0)$ 而原书印作 $(\lambda,0)$（对照本段开头的最高权向量及其权）；「对任意 $\mu\in S(\nu)$」中的 $\mu$ 应为权 $\hat\mu=(\mu,n)$，因为 $S(\nu)$ 的元素是 $\mathfrak{h}^*\oplus\C L_0$ 中的权，不是 $\mathfrak{h}^*$ 中的元素；末段「若 $\mu=\gamma(\lambda+\rho)-\rho$」中的 $\gamma$ 应为 $w\in W$（命题中 $\nu=w(\lambda+\rho)-\rho$）。[来源：本书 p.286；对照 p.285 的命题陈述与 p.286 式 (9.5-5)。]}
\end{proof}

% [Tutor Note: p.286 的最高权印作 $(\lambda,0)$，而同页上文刚说最高权是 $w(\lambda+\rho)-\rho$；此处应为 $(w(\lambda+\rho)-\rho,0)$。又 p.286 末句“若 $\mu=\gamma(\lambda+\rho)-\rho$”中的 $\gamma$ 疑为 $w$ 之误（同段其余全用 $w$）。两处均照原文，译文未改。]

\subsection{Verma 模自同态的描述}

现在用命题 9.5.1 描述 Verma 模 $\mathbb{M}_\nu$ 的自同态代数 $\operatorname{End}_{\gghat_{\kappa_c}}\mathbb{M}_\nu$。注意 $\mathbb{M}_\nu$ 由其最高权向量自由生成，故自同态由最高权向量的像唯一确定。最高权向量的像可以是任意的 $\hat{\mathfrak{n}}_+$-不变向量，且其关于 Cartan 子代数的权为 $\nu$。把这类向量所成的空间记作 $(\mathbb{M}_\nu)^{\hat{\mathfrak{n}}_+}_\nu$。用与同构 (3.3-3) 证明中相同的论证，得到
\begin{equation}
    \operatorname{End}_{\gghat_{\kappa_c}}\mathbb{M}_\nu\simeq(\mathbb{M}_\nu)^{\hat{\mathfrak{n}}_+}_\nu .
    \tag{9.5-7}
\end{equation}
于是对满足命题 9.5.1 条件的任意 $\lambda$，得到如下向量空间同构：
\begin{equation}
    \operatorname{End}_{\gghat_{\kappa_c}}\mathbb{M}_{w(\lambda+\rho)-\rho}
    \simeq(\mathbb{M}_{w(\lambda+\rho)-\rho})^{\hat{\mathfrak{n}}_+}_{w(\lambda+\rho)-\rho}
    \simeq(W^w_{w(\lambda+\rho)-\rho})^{\hat{\mathfrak{n}}_+}_{w(\lambda+\rho)-\rho}.
    \tag{9.5-8}
\end{equation}

\begin{lemma}
    对任意 $\lambda\in\mathfrak{h}^*$，$W^w_{w(\lambda+\rho)-\rho}$ 中权为 $w(\lambda+\rho)-\rho$ 的 $\hat{\mathfrak{n}}_+$-不变向量所成空间，等于其子空间 $|0\rangle_w\otimes\operatorname{Fun}\operatorname{Conn}(\Omega^{-\rho})^\lambda_D$。
\end{lemma}

\begin{proof}
    由定理 6.1.6 中 $\gghat_{\kappa_c}$ 在 Wakimoto 模上作用的公式可知，子空间
    \[
        |0\rangle_w\otimes\operatorname{Fun}\operatorname{Conn}(\Omega^{-\rho})^{w(\lambda+\rho)-\rho}_D
    \]
    中的所有向量都被 $\hat{\mathfrak{n}}_+$ 湮灭，且权为 $w(\lambda+\rho)-\rho$。下面证明 $W^w_{w(\lambda+\rho)-\rho}$ 中没有此权的其他 $\hat{\mathfrak{n}}_+$-不变向量。

    $\hat{\mathfrak{n}}_+$-不变向量特别地被李代数
    \[
        L_+\mathfrak{n}^w_-=(w\mathfrak{n}_+w^{-1}\cap\mathfrak{n}_-)\otimes t\C[t]
        \oplus(w\mathfrak{n}_+w^{-1}\cap\mathfrak{n}_+)\otimes\C[t]
    \]
    湮灭。但由定理 6.1.6 的公式可知，该李代数余自由地作用在 $M^w_{\mathfrak{g}}\otimes\operatorname{Fun}\operatorname{Conn}(\Omega^{-\rho})^\lambda_D$ 上；因此 $L_+\mathfrak{n}^w_-$-不变向量所成空间的特征标

    等于 $W^w_{w(\lambda+\rho)-\rho}$ 的特征标与 $U(L_+\mathfrak{n}^w_-)$ 的限制对偶的特征标之比。该比值等于
    \[
        e^{w(\lambda+\rho)-\rho}\prod_{n>0}(1-e^{-n\delta})^{-\ell}
        \prod_{\alpha\in w(\Delta_-)\cap\Delta_-,n\ge0}(1-e^{-n\delta-\alpha})^{-1}
        \prod_{\alpha\in w(\Delta_-)\cap\Delta_+,n>0}(1-e^{-n\delta+\alpha})^{-1}.
    \]
    因此 $M^w_{\mathfrak{g}}\otimes\operatorname{Fun}\operatorname{Conn}(\Omega^{-\rho})^\lambda_D$ 中 $L_+\mathfrak{n}^w_-$-不变向量的权 $w(\lambda+\rho)-\rho$ 子空间的特征标等于
    \[
        e^{w(\lambda+\rho)-\rho}\prod_{n>0}(1-e^{-n\delta})^{-\ell},
    \]
    这恰是子空间 $|0\rangle_w\otimes\operatorname{Fun}\operatorname{Conn}(\Omega^{-\rho})^\lambda_D$ 的特征标。证毕。
\end{proof}

% [Tutor Note: p.287 引理 9.5.2 的陈述用上标 $\lambda$，而证明第一句用上标 $w(\lambda+\rho)-\rho$；两者是同一子空间在不同参数记法下的写法（引理陈述中 $\lambda$ 是模 $W^w_\lambda$ 的参数）。引用时按引理陈述的 $\lambda$ 理解即可，勿误以为矛盾。]

回忆定理 8.3.3：中心 $Z(\gghat)$ 在 $W^w_{w(\lambda+\rho)-\rho}$ 上的作用经由同态
\[
    Z(\gghat)\simeq\operatorname{Fun}\Op_{{}^LG}(D^\times)\to\operatorname{Fun}\operatorname{Conn}(\Omega^{-\rho})_{D^\times}
\]
分解，该同态由 Miura 变换诱导。因此它经由同态
\[
    \operatorname{Fun}\Op_{{}^LG}(D^\times)\to\operatorname{Fun}\operatorname{Conn}(\Omega^{-\rho})^\lambda_D
\]
分解。但我们由 §9.3.1 的讨论知道，此同态经由同态
\[
    \operatorname{Fun}\Op^{\mathrm{RS}}_{{}^LG}(D)_{\varpi(-\lambda-\rho)}\to\operatorname{Fun}\operatorname{Conn}(\Omega^{-\rho})^\lambda_D
\]
分解。设 $\lambda$ 满足命题 9.5.1 的条件。则由命题 9.3.1，此映射是同构。\footnote{注意：为应用命题 9.3.1 的结论，需把 $G$ 换成 ${}^LG$、把 $\Omega^{\check\rho}$ 换成 $\Omega^{-\rho}$，这会改变联络留数的符号。}把它与同构 (9.5-8) 及引理 9.5.2 的结论结合，便得到如下结果（见 [Frenkel and Gaitsgory 2006c], Corollary 13.3.2）。\transnote{这一段是「自同态代数 $\simeq$ 带正则奇点 opers 的函数代数」的机制所在，分三步：定理 8.3.3 把 $Z(\gghat)$ 的作用经 $\operatorname{Fun}\Op_{{}^LG}(D^\times)\to\operatorname{Fun}\operatorname{Conn}(\Omega^{-\rho})_{D^\times}$ 分解，§9.3.1 再把它限制到 $\operatorname{Fun}\operatorname{Conn}(\Omega^{-\rho})^\lambda_D$，最后由命题 9.3.1（Miura 变换在正则奇点层上是同构）收紧到 $\operatorname{Fun}\Op^{\mathrm{RS}}_{{}^LG}(D)_{\varpi(-\lambda-\rho)}$。第三步须把 $G$ 换成 ${}^LG$、$\Omega^{\check\rho}$ 换成 $\Omega^{-\rho}$，代价是留数变号，这正是脚注所提醒的一点，也是定理 9.5.3 中指标写成 $-\nu-\rho$ 的来源。落点：§9.5.3，原书 p.287–288。[来源：本书 p.287–288；定理 8.3.3；命题 9.3.1。]}

\begin{theorem}
    中心 $Z(\gghat)$ 满射到 $\operatorname{End}_{\gghat_{\kappa_c}}\mathbb{M}_\nu$ 上，且有交换图
    \[
        \begin{CD}
            Z(\gghat) @>>> \operatorname{Fun}\Op_{{}^LG}(D^\times)\\
            @VVV @VVV\\
            \operatorname{End}_{\gghat_{\kappa_c}}\mathbb{M}_\nu @>>>
                \operatorname{Fun}\Op^{\mathrm{RS}}_{{}^LG}(D)_{\varpi(-\nu-\rho)}
        \end{CD}
    \]
    对任意权 $\nu$ 成立。此外，Verma 模 $\mathbb{M}_\nu$ 是代数 $\operatorname{End}_{\gghat_{\kappa_c}}\mathbb{M}_\nu\simeq\operatorname{Fun}\Op^{\mathrm{RS}}_{{}^LG}(D)_{\varpi(-\nu-\rho)}$ 上的自由模。
\end{theorem}

% [Tutor Note: p.288 定理 9.5.3 的两个断言（满射 + 自由模）合起来说明 $\mathbb{M}_\nu$ 是“以中心为基环的自由模”，这正是把 Verma 模看成 $\Op^{\mathrm{RS}}_{{}^LG}(D)_{\varpi(-\nu-\rho)}$ 上向量丛的代数表述，也是 §9.6 处理 Weyl 模时比较纤维的出发点。]

\section{Weyl 模的自同态}

\begin{keypoints}
    \item \textbf{核心思想：}Weyl 模 $\mathbb{V}_\lambda$（有限维最高权模诱导所得）对应幂零 oper 的子簇 $\Op^{\check\lambda}_{{}^LG}$；证明策略是构造 $\mathbb{V}_\lambda\to W_\lambda$，再比较自同态代数与 Miura 变换像。
    \item \textbf{数学工具：}Weyl 模作为 Verma 模的商；Wakimoto 模 $W_\lambda$ 及其 $0$ 次部分 $W^0_\lambda\simeq M^*_\lambda$；用射流概形计算不变量（命题 9.6.4）；屏蔽算子 $\overline{S}^{(\lambda_i)}_i$ 与顶点算子的 Fourier 系数。
    \item \textbf{关键结论：}定理 9.6.1：$Z(\gghat)\twoheadrightarrow\operatorname{End}\mathbb{V}_\lambda$ 且 $\operatorname{End}\mathbb{V}_\lambda\simeq\operatorname{Fun}\Op^{\check\lambda}_{{}^LG}$。
    \item \textbf{记号提醒：}$\mathbb{V}_{\lambda,\kappa}$ 是 Weyl 模（blackboard 体）；$\Op^{\check\lambda}_{{}^LG}$ 的下标是留数 $\varpi(-\rho-\lambda)$ 且单值性平凡。本节同时出现 $\Op^{\mathrm{nilp},\lambda}$（不带 check）与 $\Op^{\check\lambda}$（带 check）两种写法，均照原页。
\end{keypoints}

设 $V_\lambda$ 是李代数 $\mathfrak{g}$ 的有限维表示，最高权 $\lambda$ 是 $\mathfrak{g}$ 的整优权。把 $\mathfrak{g}$ 的作用扩充为 $\mathfrak{g}\otimes\C[[t]]\oplus\mathbf{1}$ 的作用，使 $\mathfrak{g}\otimes t\C[[t]]$ 平凡地作用、$\mathbf{1}$ 作用为恒等。令 $\mathbb{V}_{\lambda,\kappa}$ 为诱导 $\gghat$-模
\[
    \mathbb{V}_{\lambda,\kappa}
    =U(\gghat)\otimes_{U(\mathfrak{g}[[t]]\oplus\C\mathbf{1})}V_\lambda .
\]
这就是水平 $\kappa$、最高权 $\lambda$ 的 Weyl 模。

以下只考虑临界水平 $\kappa=\kappa_c$ 的 Weyl 模，故去掉下标 $\kappa_c$，简记为 $\mathbb{V}_\lambda$。特别地，$\mathbb{V}_0=V_{\kappa_c}(\mathfrak{g})$ 是真空模。

注意 $\mathbb{V}_\lambda=\operatorname{Ind}(V_\lambda)$，其中 $\operatorname{Ind}$ 是 §9.4.2 引入的诱导函子。由于 $V_\lambda$ 是 $\mathfrak{g}$ 上的 Verma 模 $M_\lambda$ 的有限维商，而 $\operatorname{Ind}(M_\lambda)=\mathbb{M}_\lambda$，我们得到同态
\[
    \mathbb{M}_{\lambda,\kappa_c}\to\mathbb{V}_\lambda .
\]
由定义立即可见该同态是满射，因此 Weyl 模 $\mathbb{V}_\lambda$ 实现为 Verma 模 $\mathbb{M}_\lambda$ 的商。

\subsection{结果陈述}

回忆 §9.2.3 引入的空间 $\Op^{\check\lambda}_{{}^LG}$：具有正则奇点、留数 $\varpi(-\rho-\lambda)$、单值性平凡的 ${}^LG$-oper 空间。本节将证明如下结果，它推广了对应 $\lambda=0$ 的同构 (9.5-1)。对 $\mathfrak{g}=\mathfrak{sl}_2$ 它在 [Frenkel 1995] 中证明；对一般 $\mathfrak{g}$ 在 [Frenkel and Gaitsgory 2007c] 中证明（下面的证明与该证明类似）。它也被 Beilinson 与 Drinfeld 独立地猜测（未发表）。\transnote{$\Op^{\check\lambda}_{{}^LG}$ 与 §9.5 的 $\Op^{\mathrm{RS}}_{{}^LG}(D)_{\varpi(-\nu-\rho)}$ 只差一条：单值性平凡。留数 $\varpi(-\rho-\lambda)$ 把 Weyl 模钉在某一层上，单值性条件再把它钉到该层中一个离散子簇上，因此 $\Op^{\check\lambda}_{{}^LG}$ 比整层小得多，这正是「有限维模对应幂零 oper」的几何含义。定理 9.6.1 与定理 9.5.3 因此形状完全平行，只是右侧由整层收紧为 $\Op^{\check\lambda}_{{}^LG}$。落点：§9.6.1，原书 p.289。[来源：本书 p.289 与 §9.2.3 中 $\Op^{\check\lambda}_{{}^LG}$ 的定义；$\mathfrak{sl}_2$ 情形见 [Frenkel 1995]，一般情形见 [Frenkel and Gaitsgory 2007c]（本书 p.289 引用）。]}

\begin{theorem}
    对任意整优权 $\lambda$，中心 $Z(\gghat)$ 满射到 $\operatorname{End}_{\gghat_{\kappa_c}}\mathbb{V}_\lambda$ 上，且有如下交换图
    \[
        \begin{CD}
            Z(\gghat) @>>> \operatorname{Fun}\Op_{{}^LG}(D^\times)\\
            @VVV @VVV\\
            \operatorname{End}_{\gghat_{\kappa_c}}\mathbb{V}_\lambda @>>>
                \operatorname{Fun}\Op^{\check\lambda}_{{}^LG}
        \end{CD}
    \]
\end{theorem}

% [Tutor Note: p.289 定理 9.6.1 与定理 9.5.3 的形状完全平行，只是右侧的子簇从 $\Op^{\mathrm{RS}}_{{}^LG}(D)_{\varpi(-\nu-\rho)}$ 收紧为 $\Op^{\check\lambda}_{{}^LG}$（额外要求单值性平凡）。这正是“有限维模对应幂零 oper”的几何内容：$\mathbb{V}_\lambda$ 只看到中心里落在 $\Op^{\check\lambda}_{{}^LG}$ 上的那部分函数。]

证明的策略如下：先构造同态 $\mathbb{V}_\lambda\to W_\lambda$。于是得到映射
\begin{equation}
    \mathbb{M}_\lambda\to\mathbb{V}_\lambda\to W_\lambda,
    \tag{9.6-1}
\end{equation}
以及相应的自同态代数之间的映射
\begin{equation}
    \operatorname{End}_{\gghat_{\kappa_c}}\mathbb{M}_\lambda
    \to\operatorname{End}_{\gghat_{\kappa_c}}\mathbb{V}_\lambda
    \to\operatorname{End}_{\gghat_{\kappa_c}}W_\lambda .
    \tag{9.6-2}
\end{equation}
由定理 9.5.3 与引理 6.3.4 中所用的论证，复合映射是对应于态射
\[
    \mu_\lambda:\operatorname{Conn}(\Omega^{-\rho})^\lambda_D\to\Op^{\mathrm{RS}}_{{}^LG}(D)_{\varpi(-\lambda-\rho)}
\]
（由 Miura 变换限制得到）的函数代数同态。因此复合映射的像恰为
\[
    \operatorname{Fun}\Op^{\check\lambda}_{{}^LG}\subset\operatorname{Fun}\operatorname{Conn}(\Omega^{-\rho})^\lambda_D=\operatorname{End}_{\gghat_{\kappa_c}}W_\lambda .
\]
另一方面，我们将证明 (9.6-2) 中第二个映射是单射且其像含于 $\operatorname{Fun}\Op^{\check\lambda}_{{}^LG}$。这说明 $\operatorname{End}_{\gghat_{\kappa_c}}\mathbb{V}_\lambda$ 同构于 $\operatorname{Fun}\Op^{\check\lambda}_{{}^LG}$。

\subsection{Weyl 模与 Wakimoto 模}

现在进行证明，构造同态 $\mathbb{V}_\lambda\to W_\lambda$。回忆 Wakimoto 模 $W_\lambda$ 由公式 (9.4-5) 定义：
\begin{equation}
    W_\lambda=M_{\mathfrak{g}}\otimes\operatorname{Fun}\operatorname{Conn}(\Omega^{-\rho})^\lambda_D .
    \tag{9.6-3}
\end{equation}
为构造映射 $\mathbb{V}_\lambda\to W_\lambda$，我们注意（如此前在 §6.2.3 与 §9.4.1 中所做）：常数子代数 $\mathfrak{g}\subset\gghat$ 在子空间
\[
    W^0_\lambda=\C[a^*_{\alpha,0}]_{\alpha\in\Delta_+}|0\rangle\otimes1
    \subset M_{\mathfrak{g}}\otimes\operatorname{Fun}\operatorname{Conn}(\Omega^{-\rho})^\lambda_D
\]
上的作用，与 $\mathfrak{g}$ 在反变 Verma 模 $M^*_\lambda$（实现为 $\operatorname{Fun}N_+=\C[y_\alpha]_{\alpha\in\Delta_+}$）上的自然作用一致，其中作代换 $a^*_{\alpha,0}\mapsto y_\alpha$（该作用的定义见 §5.2.4）。此外，李子代数 $\mathfrak{g}\otimes t\otimes\C[[t]]\subset\gghat_{\kappa_c}$ 在子空间 $W^0_\lambda$ 上作用为 0，而 $\mathbf{1}$ 作用为恒等。

因此，单射 $\mathfrak{g}$-同态 $V_\lambda\hookrightarrow M^*_\lambda$ 给出非零的 $\gghat_{\kappa_c}$-同态
\[
    \iota_\lambda:\mathbb{V}_\lambda\to W_\lambda,
\]
它把生成子空间 $V_\lambda\subset\mathbb{V}_\lambda$ 送到 $V_\lambda$ 在 $W^0_\lambda\simeq M^*_\lambda$ 中的像。

于是得到映射的序列
\begin{equation}
    \mathbb{M}_\lambda\to\mathbb{V}_\lambda\to W_\lambda .
    \tag{9.6-4}
\end{equation}

由公式 (9.5-7)、定理 9.5.3 与命题 9.2.1 有
\[
    \operatorname{End}_{\gghat_{\kappa_c}}\mathbb{M}_\lambda
    =(\mathbb{M}_\lambda)^{\hat{\mathfrak{n}}_+}_\lambda
    \simeq\operatorname{Fun}\Op^{\mathrm{RS}}_{{}^LG}(D)_{\varpi(-\lambda-\rho)}
    \simeq\operatorname{Fun}\Op^{\mathrm{nilp},\lambda}_{{}^LG}.
\]
其次，由引理 9.5.2 有
\[
    \operatorname{End}_{\gghat_{\kappa_c}}W_\lambda
    =(W_\lambda)^{\hat{\mathfrak{n}}_+}_\lambda
    =\operatorname{Fun}\operatorname{Conn}(\Omega^{-\rho})^\lambda_D .
\]
注意由公式 (9.6-3)，$\operatorname{Fun}\operatorname{Conn}(\Omega^{-\rho})^\lambda_D=|0\rangle\otimes\operatorname{Fun}\operatorname{Conn}(\Omega^{-\rho})^\lambda_D$ 自然地是 $W_\lambda$ 的子空间。

最后注意
\[
    \operatorname{End}_{\gghat_{\kappa_c}}\mathbb{V}_\lambda
    =(\mathbb{V}_\lambda\otimes V_\lambda^*)^{\mathfrak{g}[[t]]}
    =(\mathbb{V}_\lambda)^{\hat{\mathfrak{n}}_+}_\lambda .
\]
确实，$\mathbb{V}_\lambda$ 的任一自同态由生成子空间 $V_\lambda$ 的像唯一确定；该子空间因此定义了 $(\mathbb{V}_\lambda\otimes V_\lambda^*)^{\mathfrak{g}[[t]]}$ 中的一个 $\mathfrak{g}[[t]]$-不变元素，其最高权向量给出 $(\mathbb{V}_\lambda)^{\hat{\mathfrak{n}}_+}_\lambda$ 中的向量。与前面相同，可得所得的映射都是同构。

于是得到如下交换图：
\begin{equation}
    \begin{CD}
        (\mathbb{M}_\lambda)^{\hat{\mathfrak{n}}_+}_\lambda @>>>
            (\mathbb{V}_\lambda)^{\hat{\mathfrak{n}}_+}_\lambda @>>>
            (W_\lambda)^{\hat{\mathfrak{n}}_+}_\lambda\\
        @V{\simeq}VV @V{\simeq}VV @V{\simeq}VV\\
        \operatorname{End}_{\gghat_{\kappa_c}}\mathbb{M}_\lambda @>>>
            \operatorname{End}_{\gghat_{\kappa_c}}\mathbb{V}_\lambda @>>>
            \operatorname{End}_{\gghat_{\kappa_c}}W_\lambda\\
        @V{\simeq}VV @V{\simeq}VV @V{\simeq}VV\\
        \operatorname{Fun}\Op^{\mathrm{nilp},\lambda}_{{}^LG} @>>> ? @>>>
            \operatorname{Fun}\operatorname{Conn}(\Omega^{-\rho})^\lambda_D
    \end{CD}
    \tag{9.6-5}
\end{equation}
这说明有同态序列
\begin{equation}
    \operatorname{End}_{\gghat_{\kappa_c}}\mathbb{M}_\lambda
    \to\operatorname{End}_{\gghat_{\kappa_c}}\mathbb{V}_\lambda
    \to\operatorname{End}_{\gghat_{\kappa_c}}W_\lambda .
    \tag{9.6-6}
\end{equation}
此外还推出，$\operatorname{End}_{\gghat_{\kappa_c}}\mathbb{V}_\lambda$ 有一个子商同构于同态
\begin{equation}
    \operatorname{Fun}\Op^{\mathrm{nilp},\lambda}_{{}^LG}\to\operatorname{Fun}\operatorname{Conn}(\Omega^{-\rho})^\lambda_D
    \tag{9.6-7}
\end{equation}
的像，该同态由上图中的映射得到。

% [Tutor Note: p.291 图 (9.6-5) 中底行中间的“?”是作者有意留下的未知量：正是要证明的 $\operatorname{End}\mathbb{V}_\lambda$ 所对应的函数代数（最终为 $\operatorname{Fun}\Op^{\check\lambda}_{{}^LG}$）。图中所有竖箭头都标 $\simeq$，而底行左端到右端的水平映射不是同构；阅读时不要把 $\simeq$ 误读成水平箭头。]

由构造，最后这个同态适合交换图
\[
    \begin{CD}
        Z(\gghat) @>>> \operatorname{End}_{\gghat_{\kappa_c}}\mathbb{M}_\lambda
            @>>> \operatorname{End}_{\gghat_{\kappa_c}}W_\lambda\\
        @VVV @VVV @VVV\\
        \operatorname{Fun}\Op_{{}^LG}(D^\times) @>>>
            \operatorname{Fun}\Op^{\mathrm{nilp},\lambda}_{{}^LG} @>>>
            \operatorname{Fun}\operatorname{Conn}(\Omega^{-\rho})^\lambda_D
    \end{CD}
\]
而底行两个映射的复合 $\operatorname{Fun}\Op_{{}^LG}(D^\times)\to\operatorname{Fun}\operatorname{Conn}(\Omega^{-\rho})^\lambda_D$ 对应 $Z(\gghat)$ 在 $W_\lambda$ 上的自然作用。由定理 8.3.3，该映射就是 Miura 变换映射。因此同态 (9.6-7) 与 Miura 变换限制所得的同态一致。该同态经由 $\operatorname{Fun}\Op^{\mathrm{nilp},\lambda}_{{}^LG}$ 分解，这不过说明 $\operatorname{Conn}(\Omega^{-\rho})^\lambda_D\subset\operatorname{Conn}(\Omega^{-\rho})_{D^\times}$ 在 Miura 变换下的像属于 $\Op^{\mathrm{nilp},\lambda}_{{}^LG}$；这一点已由公式 (9.3-4) 与图 (9.3-5) 得知。此外，由图 (9.6-5) 可知，同态
\[
    \operatorname{End}_{\gghat_{\kappa_c}}\mathbb{V}_\lambda\to\operatorname{End}_{\gghat_{\kappa_c}}W_\lambda
\]
的像恰是 $\operatorname{Conn}(\Omega^{-\rho})^\lambda_D$ 在 Miura 映射下于 $\Op^{\mathrm{nilp},\lambda}_{{}^LG}$ 中的像上的函数代数。下面求出这个像。

\subsection{幂零 oper 与 Miura 变换}

回忆公式 (9.3-11) 定义的空间 $\operatorname{MOp}^{\check\lambda,w}_{{}^LG}\subset\operatorname{MOp}^{\check\lambda}_{{}^LG}$。考虑 $w=1$ 的情形。相应的空间 $\operatorname{MOp}^{\check\lambda,1}_{{}^LG}$ 参数化数对 $(\chi,\mathcal{F}'_{{}^LB,0})$，其中 $\chi=(\mathcal{F},\nabla,\mathcal{F}_{{}^LB})\in\Op^{\mathrm{nilp},\lambda}_{{}^LG}$，而 $\mathcal{F}'_{{}^LB,0}$ 是 $\chi$ 底下的 ${}^LG$-丛 $\mathcal{F}$ 在 $0\in D$ 处纤维的一个 ${}^LB$-约化，它关于留数 $\operatorname{Res}_{\check\lambda}(\chi)$ 稳定，并与 oper 约化 $\mathcal{F}_{{}^LB,0}$ 处于一般相对位置。但由于 $\operatorname{Res}_{\check\lambda}(\chi)\in{}^L\mathfrak{n}_{\mathcal{F}_{{}^LB,0}}$，这样的约化 $\mathcal{F}'_{{}^LB,0}$ 只在 $\operatorname{Res}_{\check\lambda}(\chi)=0$ 时存在，即 $\chi\in\Op^{\check\lambda}_{{}^LG}$。对满足此性质的 oper $\chi$，这类约化所成的空间是群 ${}^LN_{\mathcal{F}_{{}^LB,0}}$ 上的torsor（它同构于旗簇 ${}^LG/{}^LB$ 的大胞腔）。\transnote{torsor（主齐性空间）在此保留英文：群 ${}^LN_{\mathcal{F}_{{}^LB,0}}$ 在该空间上自由且传递地作用，后者同构于旗簇 ${}^LG/{}^LB$ 的大胞腔。这样的约化只在 $\operatorname{Res}_{\check\lambda}(\chi)=0$ 时存在，因为约化须关于留数稳定，而留数落在 ${}^L\mathfrak{n}_{\mathcal{F}_{{}^LB,0}}$ 中，非零留数会把大胞腔推出稳定化子。于是遗忘态射的像自动落在 $\Op^{\check\lambda}_{{}^LG}\subset\Op^{\mathrm{nilp},\lambda}_{{}^LG}$ 上，$\operatorname{MOp}^{\check\lambda,1}_{{}^LG}$ 成为 $\Op^{\check\lambda}_{{}^LG}$ 上的主 ${}^LN_{\mathcal{F}_{{}^LB,0}}$-丛（引理 9.6.2）。落点：§9.6.3，原书 p.292。[来源：本书 p.292；引理 9.6.2。]}

于是，遗忘态射 $\operatorname{MOp}^{\check\lambda,1}_{{}^LG}\to\Op^{\mathrm{nilp},\lambda}_{{}^LG}$ 的像是 $\Op^{\check\lambda}_{{}^LG}\subset\Op^{\mathrm{nilp},\lambda}_{{}^LG}$，且 $\operatorname{MOp}^{\check\lambda,1}_{{}^LG}$ 是 $\Op^{\check\lambda}_{{}^LG}$ 上的主 ${}^LN_{\mathcal{F}_{{}^LB,0}}$-丛。

另一方面，由定理 9.3.6 有同构 $\operatorname{Conn}(\Omega^{-\rho})^\lambda_D\simeq\operatorname{MOp}^{\check\lambda,1}_{{}^LG}$，且 Miura 变换 $\operatorname{Conn}(\Omega^{-\rho})^\lambda_D\to\Op^{\mathrm{nilp},\lambda}_{{}^LG}$ 与遗忘映射 $\operatorname{MOp}^{\check\lambda,1}_{{}^LG}\to\Op^{\mathrm{nilp},\lambda}_{{}^LG}$ 一致。于是得到：

\begin{lemma}
    $\operatorname{Conn}(\Omega^{-\rho})^\lambda_D$ 在 Miura 变换下于 $\Op^{\mathrm{nilp},\lambda}_{{}^LG}$ 中的像等于 $\Op^{\check\lambda}_{{}^LG}$。映射 $\operatorname{Conn}(\Omega^{-\rho})^\lambda_D\to\Op^{\check\lambda}_{{}^LG}$ 是 $\Op^{\check\lambda}_{{}^LG}$ 上的主 ${}^LN_{\mathcal{F}_{{}^LB,0}}$-丛。
\end{lemma}

设 $(\mathbb{V}_\lambda)^{\hat{\mathfrak{n}}_+,0}_\lambda$ 为 $(\mathbb{V}_\lambda)^{\hat{\mathfrak{n}}_+}_\lambda=\operatorname{End}_{\gghat_{\kappa_c}}\mathbb{V}_\lambda$ 中作为 (9.6-6) 左上水平同态的像的子空间。于是得到如下交换图：
\begin{equation}
    \begin{CD}
        (\mathbb{V}_\lambda)^{\hat{\mathfrak{n}}_+}_\lambda @>>>
            \operatorname{Fun}\operatorname{Conn}(\Omega^{-\rho})^\lambda_D\\
        @AAA @AAA\\
        (\mathbb{V}_\lambda)^{\hat{\mathfrak{n}}_+,0}_\lambda @>>>
            \operatorname{Fun}\Op^{\check\lambda}_{{}^LG}
    \end{CD}
    \tag{9.6-8}
\end{equation}
其中竖映射是单射，底水平映射是满射。现在要证明顶水平映射是单射，且其像含于右竖映射的像中。这将说明顶水平映射与左竖映射都是同构，从而得到定理 9.6.1 的陈述。

% [Tutor Note: p.293 图 (9.6-8) 的竖箭头朝上，表示 $(\mathbb{V}_\lambda)^{\hat{\mathfrak{n}}_+,0}_\lambda\hookrightarrow(\mathbb{V}_\lambda)^{\hat{\mathfrak{n}}_+}_\lambda$ 与 $\operatorname{Fun}\Op^{\check\lambda}_{{}^LG}\hookrightarrow\operatorname{Fun}\operatorname{Conn}(\Omega^{-\rho})^\lambda_D$ 都是单射（子代数包含）；底水平映射是满射这一点来自引理 9.6.2 的丛结构。]

第一步，将把 $\operatorname{Fun}\Op^{\check\lambda}_{{}^LG}\subset\operatorname{Fun}\operatorname{Conn}(\Omega^{-\rho})^\lambda_D$ 显式实现为若干屏蔽算子核的交，推广 §8.2.3 中 $\lambda=0$ 情形的描述。

由引理 9.6.2，$\operatorname{Fun}\Op^{\check\lambda}_{{}^LG}$ 实现为代数 $\operatorname{Fun}\operatorname{Conn}(\Omega^{-\rho})^\lambda_D$ 中由 ${}^LN_{\mathcal{F}_{{}^LB,0}}$-不变（等价地，${}^L\mathfrak{n}_{{}^LB,0}$-不变）函数组成的子代数。$\operatorname{Conn}(\Omega^{-\rho})^\lambda_D$ 的元素由形如
\begin{equation}
    \partial_t+\frac{\lambda}{t}+\mathbf{u}(t),
    \qquad
    \mathbf{u}(t)=\sum_{n\ge0}\mathbf{u}_nt^{-n-1}\in{}^L\mathfrak{h}[[t]]
    \tag{9.6-9}
\end{equation}
的算子表示。因此
\begin{equation}
    \operatorname{Fun}\operatorname{Conn}(\Omega^{-\rho})^\lambda_D
    =\C[u_{i,m}]_{i=1,\dots,\ell;m<0},
    \tag{9.6-10}
\end{equation}
其中 $u_{i,m}=\langle\check\alpha_i,\mathbf{u}_m\rangle$。这里及以下使用典范等同 ${}^L\mathfrak{h}=\mathfrak{h}^*$。现在显式计算 ${}^L\mathfrak{n}_{{}^LB,0}$ 的生成元在该空间上的作用，推广 §8.2.3 中的计算（对应 $\lambda=0$ 的情形）。

为此，固定 ${}^L\mathfrak{b}$ 中的 Cartan 子代数 ${}^L\mathfrak{h}$ 以及 $\mathcal{F}_{{}^LB,0}$ 的平凡化，使 ${}^L\mathfrak{n}_{\mathcal{F}_{{}^LB,0}}$ 与 ${}^L\mathfrak{n}$ 等同、${}^L\mathfrak{b}_{\mathcal{F}'_{{}^LB,0}}$ 与 ${}^L\mathfrak{b}_-$ 等同。接着，选取 ${}^L\mathfrak{n}$ 的生成元 $e_i$（$i=1,\dots,\ell$，关于 ${}^L\mathfrak{h}$ 在 ${}^L\mathfrak{n}$ 上的作用），使它们与先前选定的 $f_i$ 一起满足 ${}^L\mathfrak{g}$ 的标准关系。于是 ${}^L\mathfrak{n}_{\mathcal{F}_{{}^LB,0}}$ 在 $\operatorname{Conn}(\Omega^{-\rho})^\lambda_D$ 上的作用给出 $e_i$ 在 $\operatorname{Conn}(\Omega^{-\rho})^\lambda_D$ 上的无穷小作用。下面计算 $\operatorname{Fun}\operatorname{Conn}(\Omega^{-\rho})^\lambda_D$ 上相应的导子。

关于 $\mathcal{F}_{{}^LB,0}$ 的平凡化，把 Miura 变换作用于 (9.6-9) 所得的 oper 联络如下：\footnote{$\mathbf{u}(t)$ 前的负号来自：在作 Miura 变换之前，我们从 $\Omega^{-\rho}$ 上的联络换到了 $\Omega^{\rho}$ 上的联络。}
\begin{equation}
    \partial_t+\sum_{i=1}^{\ell}t^{\lambda_i}f_i-\mathbf{u}(t),
    \qquad
    \mathbf{u}(t)=\sum_{n\ge0}\mathbf{u}_nt^{-n-1}\in{}^L\mathfrak{h}[[t]],
    \tag{9.6-11}
\end{equation}
其中 $\lambda_i=\langle\check\alpha_i,\lambda\rangle$。

$e_i$ 在该联络算子上的无穷小规范作用由公式
\begin{equation}
    \delta\mathbf{u}(t)=-\left[x_i(t)\cdot e_i,\ \partial_t+\sum_{i=1}^{\ell}t^{\lambda_i}f_i-\mathbf{u}(t)\right]
    \tag{9.6-12}
\end{equation}
给出，其中 $x_i(t)\in\C[[t]]$ 满足 $x_i(0)=1$，且 (9.6-12) 的右端属于 ${}^L\mathfrak{h}[[t]]$。这些条件唯一确定 $x_i(t)$。事实上，(9.6-12) 的右端等于
\[
    -t^{\lambda_i}x_i(t)\cdot\alpha_i-u_i(t)x_i(t)\cdot e_i+\partial_tx_i(t)\cdot e_i,
\]
其中 $u_i(t)=\langle\check\alpha_i,\mathbf{u}(t)\rangle$。因此它属于 ${}^L\mathfrak{h}[[t]]$ 当且仅当
\begin{equation}
    \partial_tx_i(t)=u_i(t)x_i(t).
    \tag{9.6-13}
\end{equation}

% [Tutor Note: p.294 式 (9.6-12) 的负号与交换子次序最易出错。核对办法：右端必须落在 ${}^L\mathfrak{h}[[t]]$ 中，故 $e_i$ 的系数必须抵消——这要求 $\partial_tx_i=x_i u_i$，即 (9.6-13)；若把 (9.6-12) 中的负号丢掉，就会得到 $\partial_tx_i=-u_ix_i$，与 (9.6-14) 的指数公式差一个符号。]

若写
\[
    x_i(t)=\sum_{n\le0}x_{i,n}t^{-n},
\]
并代入公式 (9.6-13)，便得到系数 $x_{i,n}$ 满足如下递推关系：
\[
    n x_{i,n}=-\sum_{k+m=n;\,k<0;\,m\le0}u_{i,k}x_{i,m},
    \qquad n<0.
\]
由该公式与条件 $x_i(0)=1$ 得到
\begin{equation}
    \sum_{n\le0}x_{i,n}t^{-n}
    =\exp\left(\sum_{m>0}\frac{u_{i,-m}}{m}t^m\right).
    \tag{9.6-14}
\end{equation}
于是
\[
    \delta\mathbf{u}(t)=-t^{\lambda_i}x_i(t)\cdot\alpha_i,
\]
从而
\[
    \delta u_j(t)=-t^{\lambda_i}a_{ji}x_i(t),
\]
其中 $a_{ji}=\langle\check\alpha_j,\alpha_i\rangle$ 是 $\mathfrak{g}$ 的 Cartan 矩阵的 $(j,i)$ 元。换言之，算子 $e_i$ 通过导子
\begin{equation}
    \widetilde{V}_i[\lambda_i+1]
    =-\sum_{j=1}^{\ell}a_{ji}\sum_{n\ge\lambda_i}x_{i,n-\lambda_i}
    \frac{\partial}{\partial u_{j,-n-1}}
    \tag{9.6-15}
\end{equation}
作用在代数 $\operatorname{Fun}\operatorname{Conn}(\Omega^{-\rho})^\lambda_D=\C[u_{i,n}]$ 上，其中 $x_{i,n}$ 由公式 (9.6-14) 给出。

注意当 $\lambda_i=0$（$i=1,\dots,\ell$）时，这个算子出现在公式 (8.2-10) 中（相差代换 $u_{i,n}\mapsto-u_{i,n}$、$a_{ji}\mapsto a_{ij}$，因为那里考虑的是 $G$-oper 与 $\Omega^{\check\rho}$ 上的联络），并被等同于 Heisenberg 顶点代数两个 Fock 表示之间某顶点算子的 $(-1)$ 次 Fourier 系数的极限（见 §8.1.2）。我们将在命题 9.6.7 的证明中说明：算子 (9.6-15) 可以等同于同一顶点算子的 $-\!(\lambda_i+1)$ 次 Fourier 系数的极限。\transnote{指标 $-(\lambda_i+1)$ 与 (7.3-18) 中 $V_i[\lambda_i+1]$ 的指标相同，来源是 (9.6-11) 中的 $t^{\lambda_i}$ 项：$\lambda_i=\langle\check\alpha_i,\lambda\rangle$ 是整数，$t^{\lambda_i}x_i(t)$ 的最低幂次便决定 Fourier 指标的平移量 $\lambda_i+1$。$\lambda_i=0$ 时 (9.6-15) 退化为 (8.2-10)，但两处相差 $u_{i,n}\mapsto-u_{i,n}$ 与 $a_{ji}\mapsto a_{ij}$ 两个代换，因为 §8.2 在 $G$-侧、用 $\Omega^{\check\rho}$，本节在 ${}^LG$-侧、用 $\Omega^{-\rho}$；核对指标时须同时带上这两个代换。落点：§9.6.4，原书 p.295。[来源：本书 p.295 与式 (8.2-10)、(7.3-17)、(7.3-18)。]}

现在回忆，我们把 $\operatorname{Fun}\Op^{\check\lambda}_{{}^LG}$ 等同于 $\operatorname{Conn}(\Omega^{-\rho})^\lambda_D$ 上 ${}^L\mathfrak{n}$-不变函数的代数。这些恰是被 ${}^L\mathfrak{n}$ 的生成元 $e_i$（$i=1,\dots,\ell$，由公式 (9.6-15) 给出）湮灭的函数。于是得到 $\operatorname{Fun}\Op^{\check\lambda}_{{}^LG}$ 作为 $\operatorname{Fun}\operatorname{Conn}(\Omega^{-\rho})^\lambda_D$ 的子代数的如下刻画。

\begin{proposition}
    $\operatorname{Fun}\Op^{\check\lambda}_{{}^LG}$ 在 Miura 映射下于 $\operatorname{Fun}\operatorname{Conn}(\Omega^{-\rho})^\lambda_D$ 中的像，等于公式 (9.6-15) 所给出的算子 $\widetilde{V}_i[\lambda_i+1]$（$i=1,\dots,\ell$）的核之交。
\end{proposition}

作为该结果的推论，得到 $\operatorname{Fun}\Op^{\check\lambda}_{{}^LG}$ 关于算子 $L_0=-t\partial_t$ 所定义的 $\Z_+$-分次的特征标公式。\transnote{以下两个乘积与原书式 (9.6-16) 的下标都印作 $\check\alpha\in\hat\Delta_+$，但此处 $\operatorname{Fun}{}^LN$ 的生成元来自有限维 ${}^L\mathfrak{n}$，其正根集是 $\check\alpha\in\Delta_+$（紧接的正文正写作 $\check\alpha\in\Delta_+$；$\hat\Delta_+$ 是 $\gghat$ 的正根集，见 p.283 式 (9.5-4) 的用法），p.302 的式 (9.6-22) 同此。照原文保留，按 $\Delta_+$ 理解。[来源：本书 p.296、p.302；并见 p.283 式 (9.5-4)。]}由引理 9.6.2，由算子 (9.6-15) 生成的 ${}^LN$ 在 $\operatorname{Conn}(\Omega^{-\rho})^\lambda_D$ 上的作用是自由的。因此有向量空间同构
\[
    \operatorname{Fun}\operatorname{Conn}(\Omega^{-\rho})^\lambda_D
    \simeq\operatorname{Fun}\Op^{\check\lambda}_{{}^LG}\otimes\operatorname{Fun}{}^LN .
\]

由公式 (9.6-15)，算子 $e_i$ 关于上述 $\Z_+$-分次是次数为 $-(\lambda_i+1)=-\langle\check\alpha_i,\lambda+\rho\rangle$ 的齐次算子。因此 ${}^L\mathfrak{n}$ 的根生成元 $e_{\check\alpha}$（$\check\alpha\in\Delta_+$）的次数等于 $-\langle\check\alpha,\lambda+\rho\rangle$。于是 $\operatorname{Fun}{}^LN$ 的特征标等于
\[
    \prod_{\check\alpha\in\hat\Delta_+}\left(1-q^{\langle\check\alpha,\lambda+\rho\rangle}\right)^{-1}.
\]
另一方面，因 $\deg u_{i,n}=-n$，由公式 (9.6-10) 知 $\operatorname{Fun}\operatorname{Conn}(\Omega^{-\rho})^\lambda_D$ 的特征标等于
\[
    \prod_{n>0}(1-q^n)^{-\ell}.
\]
因此 $\operatorname{Fun}\Op^{\check\lambda}_{{}^LG}$ 的特征标是二者之比，把它改写成
\begin{equation}
    \prod_{\check\alpha\in\hat\Delta_+}\frac{1-q^{\langle\check\alpha,\lambda+\rho\rangle}}{1-q^{\langle\check\alpha,\rho\rangle}}
    \prod_{i=1}^{\ell}\prod_{n_i\ge d_i+1}(1-q^{n_i})^{-1},
    \tag{9.6-16}
\end{equation}
其中用到恒等式
\[
    \prod_{\check\alpha\in\Delta_+}\left(1-q^{\langle\check\alpha,\rho\rangle}\right)
    =\prod_{i=1}^{\ell}\prod_{m_i=1}^{d_i}(1-q^{m_i}).
\]

% [Tutor Note: p.296 的特征标乘积中下标印作 $\check\alpha\in\hat\Delta_+$，但这里 $\operatorname{Fun}{}^LN$ 的根生成元来自有限维 ${}^L\mathfrak{n}$，正根集应是 $\Delta_+$（同页正文正作 $\check\alpha\in\Delta_+$）；$\hat\Delta_+$ 是 $\gghat$ 的正根集。照原文保留，按 $\Delta_+$ 理解。]

\subsection{射流概形上的计算}

为完成定理 9.6.1 的证明，只需说明
\[
    \operatorname{End}_{\gghat_{\kappa_c}}\mathbb{V}_\lambda
    =(\mathbb{V}_\lambda\otimes V_\lambda^*)^{\mathfrak{g}[[t]]}
\]
的特征标不超过公式 (9.6-16) 给出的 $\operatorname{Fun}\Op^{\check\lambda}_{{}^LG}$ 的特征标。

在 $\lambda=0$ 的情形，我们在定理 8.1.3 中证明了自然映射
$\operatorname{gr}(\mathbb{V}_0)^{\mathfrak{g}[[t]]}\to(\operatorname{gr}\mathbb{V}_0)^{\mathfrak{g}[[t]]}$ 是同构。这使下述猜测显得合理：映射
\begin{equation}
    \operatorname{gr}(\mathbb{V}_\lambda\otimes V_\lambda^*)^{\mathfrak{g}[[t]]}
    \to\left(\operatorname{gr}(\mathbb{V}_\lambda\otimes V_\lambda^*)\right)^{\mathfrak{g}[[t]]}
    \tag{9.6-17}
\end{equation}
对一切 $\lambda$ 都是同构。如果我们能证明 $\left(\operatorname{gr}(\mathbb{V}_\lambda\otimes V_\lambda^*)\right)^{\mathfrak{g}[[t]]}$ 的特征标不超过 $\operatorname{Fun}\Op^{\check\lambda}_{{}^LG}$ 的特征标，那么这既说明映射 (9.6-17) 是同构，也证明了定理 9.6.1，于是大功告成。

本节与下一节将计算 $\left(\operatorname{gr}(\mathbb{V}_\lambda\otimes V_\lambda^*)\right)^{\mathfrak{g}[[t]]}$ 的特征标。但我们将发现，对一般的 $\lambda$，该特征标大于 $\operatorname{Fun}\Op^{\check\lambda}_{{}^LG}$ 的特征标。等号成立当且仅当 $\lambda$ 是 极小权。这意味着映射 (9.6-17) 是同构当且仅当 $\lambda$ 是极小的。所以沿这条路只能对极小的 $\lambda$ 证明定理 9.6.1。尽管如此，下面给出的结果可能有独立价值。为对一般的 $\lambda$ 证明定理 9.6.1，我们将使用 §9.6.6 中给出的另一个论证。

于是当前的任务是计算 $\left(\operatorname{gr}(\mathbb{V}_\lambda\otimes V_\lambda^*)\right)^{\mathfrak{g}[[t]]}$ 的特征标。$U_{\kappa_c}(\gghat)$ 上的 PBW 滤过给出 $\mathbb{V}_\lambda$（看作由子空间 $V_\lambda$ 生成的 $\gghat_{\kappa_c}$-模）上的滤过。相应的分次伴随空间是
\[
    \operatorname{gr}\mathbb{V}_\lambda
    =V_\lambda\otimes\Sym(\mathfrak{g}((t))/\mathfrak{g}[[t]])
    \simeq V_\lambda\otimes\operatorname{Fun}\mathfrak{g}^*[[t]].
\]
因此得到嵌入
\[
    \left(\operatorname{gr}(\mathbb{V}_\lambda\otimes V_\lambda^*)\right)^{\mathfrak{g}[[t]]}
    \hookrightarrow\left(V_\lambda\otimes V_\lambda^*\otimes\operatorname{Fun}\mathfrak{g}^*[[t]]\right)^{\mathfrak{g}[[t]]}.
\]
于是可通过计算 $\left(V_\lambda\otimes V_\lambda^*\otimes\operatorname{Fun}\mathfrak{g}^*[[t]]\right)^{\mathfrak{g}[[t]]}$ 的特征标，得到 $\operatorname{End}_{\gghat_{\kappa_c}}\mathbb{V}_\lambda$ 的特征标的一个上界。

第一步是计算 $\operatorname{Fun}\mathfrak{g}^*[[t]]$ 在 Lie 代数 $t\mathfrak{g}[[t]]\subset\mathfrak{g}[[t]]$ 作用下的不变量代数。

回忆空间
\[
    \mathcal{P}=\mathfrak{g}^*/G=\Spec(\operatorname{Fun}\mathfrak{g}^*)^G\simeq\Spec\C[\bar P_i]_{i=1,\dots,\ell},
\]
以及态射 $p:\mathfrak{g}^*\to\mathcal{P}$。考虑相应的射流概形
\[
    J_N\mathcal{P}=\Spec\C[\bar P_{i,n}]_{i=1,\dots,\ell;-N-1\le n<0},
\]
它在定理 3.4.2 的证明中引入。这里及以下允许 $N$ 有限或无穷（射流概形的定义见 §3.4.2）。

由定理 3.4.2 有
\[
    (\operatorname{Fun}J_N\mathfrak{g}^*)^{J_NG}=\C[\bar P_{i,n}]_{i=1,\dots,\ell;-N-1\le n<0}.
\]
换言之，态射 $J_Np:J_N\mathfrak{g}^*\to J_N\mathcal{P}$ 经由同构
\[
    J_N\mathfrak{g}^*/J_NG\simeq J_N\mathcal{P}
\]
分解，其中按定义
\[
    J_N\mathfrak{g}^*/J_NG=\Spec(\operatorname{Fun}J_N\mathfrak{g}^*)^{J_NG}.
\]
我们有自然映射 $\mathfrak{g}^*\to\mathcal{P}$ 与 $J_N\mathcal{P}\to\mathcal{P}$。把纤维积记作
\[
    \widetilde{J_N\mathcal{P}}=J_N\mathcal{P}\times_{\mathcal{P}}\mathfrak{g}^*.
\]
按定义，$\widetilde{J_N\mathcal{P}}$ 上的函数代数 $\operatorname{Fun}\widetilde{J_N\mathcal{P}}$ 由 $\operatorname{Fun}J_N\mathcal{P}=(\operatorname{Fun}J_N\mathfrak{g}^*)^{J_NG}$ 与 $\operatorname{Fun}\mathfrak{g}^*$ 生成，唯一的约束是把 $(\operatorname{Fun}J_N\mathfrak{g}^*)^{J_NG}$ 与 $\operatorname{Fun}\mathfrak{g}^*$ 中由 $(\operatorname{Fun}\mathfrak{g}^*)^G$ 的同一元素拉回所得的元素等同。

考虑态射
\[
    \widetilde{J_Np}:J_N\mathfrak{g}^*\to\widetilde{J_N\mathcal{P}},
    \qquad
    A(t)\mapsto(J_Np(A(t)),A(0)),
\]
以及相应的代数同态
\begin{equation}
    \operatorname{Fun}\widetilde{J_N\mathcal{P}}\to\operatorname{Fun}J_N\mathfrak{g}^* .
    \tag{9.6-18}
\end{equation}
设 $J_NG^{(1)}$ 是 $J_NG$ 中模 $t$ 与恒等合一的元素所成的子群。注意 $J_NG$ 的李代数是 $J_N\mathfrak{g}=\mathfrak{g}\otimes(\C[[t]]/t^{N+1}\C[[t]])$，而 $J_NG^{(1)}$ 的李代数是 $J_N\mathfrak{g}^{(1)}=\mathfrak{g}\otimes(t\C[[t]]/t^{N+1}\C[[t]])$。对本讨论中所考虑的一切代数，$J_NG$（相应地 $J_NG^{(1)}$）下的不变空间都等于 $J_N\mathfrak{g}$（相应地 $J_N\mathfrak{g}^{(1)}$）下的相应不变空间。

显然，同态 (9.6-18) 的像由 $\operatorname{Fun}J_N\mathfrak{g}^*$ 中在 $J_N\mathfrak{g}^{(1)}$（从而 $J_NG^{(1)}$）作用下不变的元素组成。

\begin{proposition}
    同态 (9.6-18) 是单射，其像等于 $\operatorname{Fun}J_N\mathfrak{g}^*$ 中 $J_N\mathfrak{g}^{(1)}$-不变元素的代数（$N$ 有限或无穷均如此）。
\end{proposition}

\begin{proof}
    与定理 3.4.2 的证明一样，考虑 $\mathfrak{g}^*$ 中的正则元素子集 $\mathfrak{g}^*_{\mathrm{reg}}$。则有光滑满射 $p_{\mathrm{reg}}:\mathfrak{g}^*\to\mathcal{P}$。相应的态射 $J_Np_{\mathrm{reg}}:J_N\mathfrak{g}^*_{\mathrm{reg}}\to J_N\mathcal{P}$ 也光滑且满。由射流概形的定义，
    \[
        J_N\mathfrak{g}^*_{\mathrm{reg}}
        \simeq\mathfrak{g}^*_{\mathrm{reg}}\times\left(\mathfrak{g}^*\otimes(t\C[[t]]/t^{N+1}\C[[t]])\right).
    \]
    因此 $J_N\mathfrak{g}^*_{\mathrm{reg}}$ 是 $J_N\mathfrak{g}^*$ 的开稠密光滑子概形，其补的余维大于 1。

    如定理 3.4.2 的证明所说明，$J_Np_{\mathrm{reg}}$ 的每条纤维是群 $J_NG$ 的单个轨道。我们曾用这一事实推出
    \[
        (\operatorname{Fun}J_N\mathfrak{g}^*)^{J_NG}
        =(\operatorname{Fun}J_N\mathfrak{g}^*_{\mathrm{reg}})^{J_NG}\simeq J_N\mathcal{P}.
    \]
    现在要考虑的不是 $J_NG$ 下的不变量，而是子群 $J_NG^{(1)}\subset J_NG$ 下的不变量。同样有
    \[
        (\operatorname{Fun}J_N\mathfrak{g}^*)^{J_NG^{(1)}}
        =(\operatorname{Fun}J_N\mathfrak{g}^*_{\mathrm{reg}})^{J_NG^{(1)}},
    \]
    故集中考虑后者。

    先设 $N$ 有限。先证 (9.6-18) 是单射。令
    \[
        \widetilde{J_N\mathcal{P}}_{\mathrm{reg}}=J_N\mathcal{P}\times_{\mathcal{P}}\mathfrak{g}^*_{\mathrm{reg}}.
    \]
    考虑态射
    \[
        \widetilde{J_Np}_{\mathrm{reg}}:J_N\mathfrak{g}^*_{\mathrm{reg}}\to\widetilde{J_N\mathcal{P}}_{\mathrm{reg}},
        \qquad
        A(t)\mapsto(J_Np_{\mathrm{reg}}(A(t)),A(0)).
    \]
    断言它满。事实上，设给定点 $(\phi,x)\in\widetilde{J_N\mathcal{P}}_{\mathrm{reg}}$。由于 $J_Np_{\mathrm{reg}}$ 满，存在 $A(t)\in J_N\mathfrak{g}^*_{\mathrm{reg}}$ 使 $J_Np_{\mathrm{reg}}(A(t))=\phi$。但此时 $A(0)$ 与 $x$ 是 $\mathfrak{g}^*_{\mathrm{reg}}$ 的元素，且 $p_{\mathrm{reg}}(A(0))=p_{\mathrm{reg}}(x)$。已知群 $G$ 沿 $p_{\mathrm{reg}}$ 的纤维传递地作用，故 $x=g\cdot A(0)$（某 $g\in G$）。于是 $\widetilde{J_Np}_{\mathrm{reg}}(g\cdot A(t))=(\phi,x)$。故 $\widetilde{J_Np}_{\mathrm{reg}}$ 满，从而 (9.6-18) 是单射。

    群 $J_NG^{(1)}$ 沿 $\widetilde{J_Np}_{\mathrm{reg}}$ 的纤维作用。断言每条纤维是单个 $J_NG^{(1)}$-轨道。

    为证明，描述 $J_N\mathfrak{g}^*_{\mathrm{reg}}$ 的点在群 $J_NG$ 作用下的稳定化子。$J_NG$ 的任意元素可写成
    \[
        g(t)=g_0\exp\left(\sum_{m=1}^{N}y_mt^m\right),
        \qquad g_0\in G,\ y_m\in\mathfrak{g}.
    \]
    $g(t)$ 稳定 $J_N\mathfrak{g}^*_{\mathrm{reg}}$ 的点
    \[
        x(t)=\sum_{n=0}^{N}x_nt^n,
        \qquad x_0\in\mathfrak{g}^*_{\mathrm{reg}};\ x_n\in\mathfrak{g}^*,\ n>0,
    \]
    的条件等价于关于 $g_0$ 与 $y_1,\dots,y_m$ 的递推关系组。第一个关系是 $g_0\cdot x_0=x_0$，即 $g_0$ 属于 $x_0$ 的稳定化子。下一个关系是
    \[
        \operatorname{ad}^*_{x_0}(y_1)=x_1-g_0^{-1}\cdot x_1.
    \]
    容易看出，对正则元素 $x_0$ 与任意 $x_1\in\mathfrak{g}^*$，右端总在 $\operatorname{ad}^*_{x_0}$ 的像中。因此该关系可对 $y_1$ 解出。但 $x_0$ 的中心化子维数为 $\ell=\operatorname{rank}(\mathfrak{g})$，故得到 $y_1$ 的 $\ell$ 维解族。

    归纳地继续解这些关系。第 $m$ 个递推关系形如：$\operatorname{ad}^*_{x_0}(y_m)$ 等于由前面各关系的解完全确定的 $\mathfrak{g}^*$ 中元素。若该元素在 $\operatorname{ad}^*_{x_0}$ 的像中，则第 $m$ 个关系可对 $y_m$ 解出，得到 $\ell$ 维解族。否则无解，递推过程停止。这意味着稳定化子中没有满足 $g(0)=g_0$ 的 $g(t)\in J_NG$。

    于是 $J_N\mathfrak{g}^*_{\mathrm{reg}}$ 中任意 $x(t)$ 在 $J_NG$ 中中心化子的可能最大维数等于 $(N+1)\ell$。该界被达到，当且仅当对一般的 $g_0\in\operatorname{Stab}_{x_0}(G)$ 所有递推关系都可解（即第 $m$ 个关系的右端对一切 $m=1,\dots,N$ 都属于 $\operatorname{ad}^*_{x_0}$ 的像）。但若这些关系对一般的 $g_0$ 可解，则对一切 $g_0\in\operatorname{Stab}_{x_0}(G)$ 都可解。于是得到：$x(t)\in J_N\mathfrak{g}^*_{\mathrm{reg}}$ 在 $J_NG$ 中的中心化子维数为 $(N+1)\ell$，当且仅当对任意 $g_0\in\operatorname{Stab}_{x_0}(G)$ 所有递推关系都可解。

    然而注意，$J_N\mathfrak{g}^*_{\mathrm{reg}}$ 的维数等于 $(N+1)\dim\mathfrak{g}$，$J_N\mathcal{P}$ 的维数等于 $(N+1)\ell$。因此 $J_Np_{\mathrm{reg}}$ 的一般纤维的维数为 $(N+1)(\dim\mathfrak{g}-\ell)$。由于群 $J_NG$ 沿纤维传递地作用，$J_N\mathfrak{g}^*_{\mathrm{reg}}$ 的一般点的稳定化子维数为 $(N+1)\ell$。因此对一般的 $x(t)\in J_N\mathfrak{g}^*_{\mathrm{reg}}$，上述一切递推关系对任意 $g_0\in\operatorname{Stab}_{x_0}$ 都可解。于是它们对一切 $x(t)$ 与 $g_0\in\operatorname{Stab}_{x_0}$ 都可解。特别地，对任意 $x(t)\in J_N\mathfrak{g}^*_{\mathrm{reg}}$ 与 $g_0\in\operatorname{Stab}_{x_0}(G)$，存在 $g(t)\in J_NG$ 使 $g(0)=g_0$ 且 $g(t)\cdot x(t)=x(t)$。

    现在可以证明 $\widetilde{J_Np}_{\mathrm{reg}}$ 的每条纤维是单个 $J_NG^{(1)}$-轨道。考虑同一纤维中的两点 $x_1(t),x_2(t)\in J_N\mathfrak{g}^*_{\mathrm{reg}}$。特别地 $x_1(0)=x_2(0)=x_0$。这两点也在 $J_Np_{\mathrm{reg}}$ 的同一纤维中，故存在 $g(t)\in J_NG$ 使 $x_2(t)=g(t)\cdot x_1(t)$。特别地 $g(0)\cdot x_0=x_0$。按前述讨论，存在 $\tilde{g}(t)\in J_NG$ 使 $\tilde{g}(0)=g(0)$ 且 $\tilde{g}(t)\cdot x_2(t)=x_2(t)$。于是 $\tilde{g}(t)^{-1}g(t)\in J_NG^{(1)}$ 具有所需性质：它把 $x_1(t)$ 变为 $x_2(t)$。

    于是 $\widetilde{J_Np}_{\mathrm{reg}}$ 的每条纤维是单个 $J_NG^{(1)}$-轨道。与定理 3.4.2 的证明相同，这说明 $\operatorname{Fun}J_N\mathfrak{g}^*_{\mathrm{reg}}$（从而 $\operatorname{Fun}J_N\mathfrak{g}^*$）中的 $J_NG^{(1)}$-不变元素的代数，等于 $\operatorname{Fun}\widetilde{J_N\mathcal{P}}$ 在单射同态 (9.6-18) 下的像。至此对有限的 $N$ 完成了证明。

    最后考虑无穷射流概形的情形。如前，把 $J_\infty X$ 记作 $JX$。按定义，$\operatorname{Fun}\widetilde{J\mathcal{P}}$ 的任一元素 $A$ 都来自某个足够大的 $N$ 的 $\operatorname{Fun}\widetilde{J_N\mathcal{P}}$ 中元素 $A_N$ 的拉回。$A$ 在 $\operatorname{Fun}J\mathfrak{g}^*$ 中的像是 $A_N$ 在 $\operatorname{Fun}J_N\mathfrak{g}^*$ 中的像的拉回。因此 (9.6-18) 对有限 $N$ 的单射性蕴含其对 $N=\infty$ 的单射性。

    $J\mathfrak{g}^*_{\mathrm{reg}}$ 上任一 $JG^{(1)}$-不变函数都来自某个 $J_N\mathfrak{g}^*_{\mathrm{reg}}$ 上 $J_NG^{(1)}$-不变函数的拉回。于是 $J\mathfrak{g}^*_{\mathrm{reg}}$ 上 $JG^{(1)}$-不变函数的代数等于 $\operatorname{Fun}\widetilde{J\mathcal{P}}$。
\end{proof}

% [Tutor Note: p.298–p.301 命题 9.6.4 的证明结构与定理 3.4.2 平行，但把“$J_NG$-不变量”换成“$J_NG^{(1)}$-不变量”。关键两步：(1) 用纤维积 $\widetilde{J_N\mathcal{P}}_{\mathrm{reg}}=J_N\mathcal{P}\times_{\mathcal{P}}\mathfrak{g}^*_{\mathrm{reg}}$ 把 $J_Np_{\mathrm{reg}}$ 的满性“拼上 $A(0)$”；(2) 用递推关系的可解性把中心化子维数算到 $(N+1)\ell$，从而证明每条纤维是单个 $J_NG^{(1)}$-轨道。此处“余维大于 1 ⇒ 不变函数由开子集上的值唯一决定”是定理 3.4.2 已用过的同一手法。]

下面这个命题的直截了当的推广是有用的。对如上 $N$，设 $k$ 是小于 $N$ 的非负整数。把纤维积记作
\[
    \widetilde{J_N^{(k)}\mathcal{P}}=J_N\mathcal{P}\times_{J_k\mathcal{P}}J_k\mathfrak{g}^*,
\]
其中取自然态射 $J_N\mathcal{P}\to J_k\mathcal{P}$ 与 $J_k\mathfrak{g}^*\to J_k\mathcal{P}$。

考虑态射
\[
    \widetilde{J_N^{(k)}p}:J_N\mathfrak{g}^*\to\widetilde{J_N^{(k)}\mathcal{P}},
    \qquad
    A(t)\mapsto\left(J_Np(A(t)),A(t)\ \mathrm{mod}\ (t^{k+1})\right),
\]
以及相应的代数同态
\begin{equation}
    \operatorname{Fun}\widetilde{J_N^{(k)}\mathcal{P}}\to\operatorname{Fun}J_N\mathfrak{g}^* .
    \tag{9.6-19}
\end{equation}
设 $J_NG^{(k+1)}$ 是 $J_NG$ 中模 $(t^{k+1})$ 与恒等合一的元素所成的子群。注意 $J_NG^{(k+1)}$ 的李代数是
\[
    J_N\mathfrak{g}^{(k+1)}=\mathfrak{g}\otimes\left(t^{k+1}\C[[t]]/t^{N+1}\C[[t]]\right).
\]
重复命题 9.6.4 证明中的论证，得到命题 9.6.4 的如下推广（后者对应 $k=0$ 的特殊情形）：

\begin{proposition}
    同态 (9.6-19) 是单射，其像等于 $\operatorname{Fun}J_N\mathfrak{g}^*$ 中 $J_N\mathfrak{g}^{(k+1)}$-不变元素的代数（$N$ 有限或无穷均如此）。
\end{proposition}

% [Tutor Note: p.301 的推广命题对应“只要求模 $t^{k+1}$ 与恒等合一”的子群 $J_NG^{(k+1)}$，其用途在 §9.6.5：取 $k=0$ 得 $JG^{(1)}$-不变量，再逐级过滤可把中心自同态分解为射流概形上的各级不变量。注意 $\widetilde{J_N^{(k)}\mathcal{P}}$ 的下标是 $J_k\mathcal{P}$（两次）——这是射流概形的逐级相容性。]

\subsection{分次伴随空间中的不变子空间}

群 $G$ 与 Lie 代数 $\mathfrak{g}$ 自然地作用在 $J\mathfrak{g}^*=\mathfrak{g}^*[[t]]$ 上 $JG^{(1)}$-不变多项式函数所成的空间上，并保持由多项式次数定义的 $\operatorname{Fun}\mathfrak{g}^*[[t]]$ 上的自然 $\Z_+$-分次。由命题 9.6.4，这个 $JG^{(1)}$-不变函数空间同构于 $\operatorname{Fun}\widetilde{J\mathcal{P}}$。作为 $\Z_+$-分次 $\mathfrak{g}$-模，后者

同构于
\[
    \operatorname{Fun}\mathfrak{g}^*/(\operatorname{Inv}\mathfrak{g}^*)_+
    \otimes_{\C}\C[\bar P_{i,n}]_{i=1,\dots,\ell;n<0}.
\]
这里 $(\operatorname{Inv}\mathfrak{g}^*)_+$ 表示 $\operatorname{Fun}\mathfrak{g}^*$ 中由 $\mathfrak{g}$-不变子代数 $\operatorname{Inv}\mathfrak{g}^*$ 的增广理想生成的理想。$\mathfrak{g}$ 的作用只在前一个因子上非平凡，$\Z_+$-分次由两个因子上的分次合并得到。

因此有 $\Z_+$-分次向量空间的同构
\begin{equation}
    \begin{aligned}
        \left(V_\lambda\otimes V_\lambda^*\otimes\operatorname{Fun}\mathfrak{g}^*[[t]]\right)^{\mathfrak{g}[[t]]}
        =\;&\operatorname{Hom}_{\mathfrak{g}}\left(V_\lambda\otimes V_\lambda^*,
            \operatorname{Fun}\mathfrak{g}^*/(\operatorname{Inv}\mathfrak{g}^*)_+\right)\\
        &\otimes_{\C}\C[\bar P_{i,n}]_{i=1,\dots,\ell;n<0}.
    \end{aligned}
    \tag{9.6-20}
\end{equation}
下面计算 (9.6-20) 的特征标。

由 B. Kostant 的一个定理 [1963]，对任意 $\mathfrak{g}$-模 $V$，$\Z_+$-分次空间
$\operatorname{Hom}_{\mathfrak{g}}(V,\operatorname{Fun}\mathfrak{g}^*/(\operatorname{Inv}\mathfrak{g}^*)_+)$ 同构于 $V$ 的 $\mathfrak{g}_e$-不变子空间，其中 $e$ 是 $\mathfrak{g}$ 的正则幂零元，$\mathfrak{g}_e$ 是它的中心化子（$\mathfrak{g}$ 的 $\ell$ 维交换李子代数）。不失一般性，取 $e$ 为元素 $\sum_i e_i$，其中 $\{e_i\}_{i=1,\dots,\ell}$ 是 $\mathfrak{n}$ 的生成元。则 $e$ 含于 $\mathfrak{sl}_2$-三元组 $\{f,\check\rho,e\}\subset\mathfrak{g}$。$V^{\mathfrak{g}_e}$ 上相应的 $\Z_+$-分次是由 $\check\rho$ 的作用（带适当平移）定义的 $\Z_+$-分次（见下）。\transnote{这里用了 B. Kostant 的定理 [Kostant 1963]：$\operatorname{Hom}_{\mathfrak{g}}(V,\operatorname{Fun}\mathfrak{g}^*/(\operatorname{Inv}\mathfrak{g}^*)_+)\simeq V^{\mathfrak{g}_e}$，其中 $\mathfrak{g}_e$ 是正则幂零元 $e$ 的中心化子。取 $e=\sum_i e_i$ 后 $e$ 与 $f$、$\check\rho$ 组成主 $\mathfrak{sl}_2$-三元组，$V^{\mathfrak{g}_e}$ 上的 $\Z_+$-分次由 $\check\rho$ 的本征值给出；约定 $e$ 的次数为 $1$、$f$ 的次数为 $-1$ 时，平移量恰是 $-\langle\lambda,\check\rho\rangle$（因最低权向量 $u_\lambda$ 的权为 $-\langle\lambda,\check\rho\rangle$）。若把三元组写成 $\{e,\check\rho,f\}$ 而误认 $e,f$ 的位置，整体符号会反过来。落点：§9.6.5，原书 p.302。[来源：本书 p.302；B. Kostant, \emph{Lie group representations on polynomial rings}, Amer. J. Math. 85 (1963), 327--404；[Kac 1990] 式 (10.9.4)（本书 p.302 引用）。]}

于是得到：(9.6-20) 右端张量积中第一个因子
\begin{equation}
    \operatorname{Hom}_{\mathfrak{g}}\left(V_\lambda\otimes V_\lambda^*,
        \operatorname{Fun}\mathfrak{g}^*/(\operatorname{Inv}\mathfrak{g}^*)_+\right)
    \tag{9.6-21}
\end{equation}
的特征标等于
\[
    (V_\lambda\otimes V_\lambda^*)^{\mathfrak{g}_e}
    =\operatorname{Hom}_{U(\mathfrak{g}_e)}(V_\lambda,V_\lambda).
\]
若 $V_\lambda$ 是多项式代数 $U(\mathfrak{g}_e)$ 上的循环模、由最低权向量 $u_\lambda$ 生成，则 $U(\mathfrak{g}_e)$-模同态 $V_\lambda\to V_\lambda$ 由 $u_\lambda$ 的像唯一确定。于是得到
\[
    \operatorname{Hom}_{U(\mathfrak{g}_e)}(V_\lambda,V_\lambda)\simeq V_\lambda,
\]
其中右端的 $\Z$-分次是 $V_\lambda$ 上的主分次减去 $\langle\lambda,\check\rho\rangle$（注意 $u_\lambda$ 的权等于 $-\langle\lambda,\check\rho\rangle$）。于是 (9.6-21) 的特征标等于 $V_\lambda$ 的主特征标乘以 $q^{-\langle\lambda,\check\rho\rangle}$。由 [Kac 1990] 的公式 (10.9.4)，后者等于
\begin{equation}
    \operatorname{ch}V_\lambda
    =\prod_{\check\alpha\in\hat\Delta_+}\frac{1-q^{\langle\check\alpha,\lambda+\rho\rangle}}{1-q^{\langle\check\alpha,\rho\rangle}}.
    \tag{9.6-22}
\end{equation}

% [Tutor Note: p.302 的 Kostant 定理把 $\operatorname{Hom}_{\mathfrak{g}}(V,\operatorname{Fun}\mathfrak{g}^*/(\operatorname{Inv}\mathfrak{g}^*)_+)$ 换成 $\mathfrak{g}_e$-不变量；这里 $e=\sum e_i$ 是正则幂零元，$\{f,\check\rho,e\}$ 是主 $\mathfrak{sl}_2$-三元组，故 $e$ 的权（在 $\check\rho$-分次下）为 1，$f$ 为 $-1$。若把三元组写成 $\{e,\check\rho,f\}$ 而弄错 $e,f$ 的位置，主分次的平移方向会整体反过来。另注意 (9.6-22) 的下标又是 $\hat\Delta_+$，同 p.296 的问题，应读作 $\Delta_+$。]

因此，所求的 $\left(V_\lambda\otimes V_\lambda^*\otimes\operatorname{Fun}\mathfrak{g}^*[[t]]\right)^{\mathfrak{g}[[t]]}$ 的特征标（等于 (9.6-20) 右端的特征标）将由公式 (9.6-16) 给出。

然而 $V_\lambda$ 是 $U(\mathfrak{g}_e)$ 上的循环模，当且仅当 $\lambda$ 是 极小权（即 $V_\lambda$ 中每个权的重数为一）。“仅当”部分来自：$(V_\lambda\otimes V_\lambda^*)^{\mathfrak{g}_e}$ 的维数等于 $V_\lambda\otimes V_\lambda^*$ 中权 0 子空间的维数（见 [Kostant 1963]），后者等于
\[
    \sum_\mu\left(\dim V_\lambda(\mu)\right)^2>\sum_\mu\dim V_\lambda(\mu)=\dim V_\lambda,
\]
除非 $\lambda$ 是极小的（这里 $V_\lambda(\mu)$ 表示 $V_\lambda$ 的权 $\mu$ 分量）。“当”部分是 V. Ginzburg 的结果，见 [1995], Prop. 1.8.1（另见 [Ginzburg 1998]）。

于是对 极小权 $\lambda$，上述计算说明 $\operatorname{End}_{\gghat_{\kappa_c}}\mathbb{V}_\lambda=(\mathbb{V}_\lambda\otimes V_\lambda^*)^{\mathfrak{g}[[t]]}$（其分次伴随含于空间 $\left(V_\lambda\otimes V_\lambda^*\otimes\operatorname{Fun}\mathfrak{g}^*[[t]]\right)^{\mathfrak{g}[[t]]}$ 中）的特征标不超过 (9.6-16) 给出的 $\operatorname{Fun}\Op^{\check\lambda}_{{}^LG}$ 的特征标。又由 §9.6.2 的结果已知 $\operatorname{Fun}\Op^{\check\lambda}_{{}^LG}$ 同构于 $\operatorname{End}_{\gghat_{\kappa_c}}\mathbb{V}_\lambda$ 的一个子商，于是得到定理 9.6.1 对极小最高权 $\lambda$ 的陈述。

但对非极小最高权 $\lambda$，我们发现 $\left(V_\lambda\otimes V_\lambda^*\otimes\operatorname{Fun}\mathfrak{g}^*[[t]]\right)^{\mathfrak{g}[[t]]}$ 的特征标严格大于 $\operatorname{Fun}\Op^{\check\lambda}_{{}^LG}$ 的特征标，故上述计算不能用来证明定理 9.6.1。（此外，这说明映射 (9.6-17) 是同构当且仅当 $\lambda$ 是极小的。）下一节给出另一个论证。\transnote{极小权（每个权的重数都是 1）恰是 $V_\lambda$ 成为 $U(\mathfrak{g}_e)$ 上循环模的情形；一旦 $\lambda$ 非极小，$\sum_\mu(\dim V_\lambda(\mu))^2>\dim V_\lambda$ 使 $(V_\lambda\otimes V_\lambda^*)^{\mathfrak{g}_e}$ 严格大于 $V_\lambda$，于是 $\left(V_\lambda\otimes V_\lambda^*\otimes\operatorname{Fun}\mathfrak{g}^*[[t]]\right)^{\mathfrak{g}[[t]]}$ 的特征标严格大于 $\operatorname{Fun}\Op^{\check\lambda}_{{}^LG}$ 的，比较法失效。括号中的断言（(9.6-17) 是同构当且仅当 $\lambda$ 为极小）同出一源；非极小情形必须改走 §9.6.6 的屏蔽算子路线。「当」的部分依赖 V. Ginzburg 的结果 [Ginzburg 1995], Prop. 1.8.1，并非初等。落点：§9.6.5，原书 p.303。[来源：本书 p.303；[Kostant 1963] 与 [Ginzburg 1995]（本书 p.303 引用）。]}

\subsection{证明的完成}

我们将通过说明图 (9.6-8) 的顶水平同态
\begin{equation}
    (\mathbb{V}_\lambda\otimes V_\lambda^*)^{\mathfrak{g}[[t]]}
    =(\mathbb{V}_\lambda)^{\hat{\mathfrak{n}}_+}_\lambda
    \to\operatorname{Fun}\operatorname{Conn}(\Omega^{-\rho})^\lambda_D
    \tag{9.6-23}
\end{equation}
是单射、且其像含于公式 (9.6-15) 所给算子 $\widetilde{V}_i[\lambda_i+1]$（$i=1,\dots,\ell$）的核之交，来完成定理 9.6.1 的证明。

如下陈述（推广命题 7.1.1，对应 $\lambda=0$ 的情形）在 [Frenkel and Gaitsgory 2007c] 中证明。

\begin{theorem}
    对任意整优权 $\lambda$，映射 $\iota_\lambda:\mathbb{V}_\lambda\to W_\lambda$ 是单射。
\end{theorem}

因此映射 (9.6-23) 是单射。由引理 9.5.2，$(W_\lambda)^{\hat{\mathfrak{n}}_+}_\lambda$ 是子空间
\[
    \operatorname{Fun}\operatorname{Conn}(\Omega^{-\rho})^\lambda_D=\pi_\lambda\subset W_\lambda .
\]
于是得到单射
\begin{equation}
    (\mathbb{V}_\lambda)^{\hat{\mathfrak{n}}_+}_\lambda\hookrightarrow\operatorname{Fun}\pi_\lambda .
    \tag{9.6-24}
\end{equation}

\begin{proposition}
    映射 (9.6-24) 的像含于公式 (9.6-15) 所给算子 $\widetilde{V}_i[\lambda_i+1]$（$i=1,\dots,\ell$）的核之交。
\end{proposition}

\begin{proof}
    使用与命题 7.3.6（对应 $\lambda=0$ 的情形）相同的证明。§7.3.5 中，对任意整优权 $\lambda$ 构造了屏蔽算子
    \[
        \overline{S}^{(\lambda_i)}_i:W_{\lambda,\kappa_c}\to\widetilde{W}_{0,\lambda,\kappa_c},
    \]
    它们与 $\gghat_{\kappa_c}$ 的作用交换。由这些算子的公式 (7.3-17) 显然可知，它们湮灭 $W_{\lambda,\kappa_c}$ 中的最高权向量。由于 $\mathbb{V}_\lambda$ 在 $W_\lambda=W_{\lambda,\kappa_c}$ 中的像由该向量生成，故该像含于算子 $\overline{S}^{(\lambda_i)}_i$ 的核之交。但 $(\mathbb{V}_\lambda)^{\hat{\mathfrak{n}}_+}_\lambda$ 的像也含于子空间 $\pi_\lambda\subset W_\lambda$，因此含于把 $\overline{S}^{(\lambda_i)}_i$ 限制到 $\pi_\lambda\subset W_\lambda$ 所得算子的核之交。如 §7.3.5 末尾所说明，该限制等于公式 (7.3-18) 所给的算子 $V_i[\lambda_i+1]$。

    比较公式 (7.3-18) 与 (9.6-15) 可知，在代换 $b_{i,n}\mapsto u_{i,n}$ 下，算子 $V_i[\lambda_i+1]$ 变成公式 (9.6-15) 所给的算子 $\widetilde{V}_i[\lambda_i+1]$。证毕。
\end{proof}

现在，映射 (9.6-23) 的单射性、命题 9.6.7 与命题 9.6.3 表明映射
\[
    (\mathbb{V}_\lambda)^{\hat{\mathfrak{n}}_+}_\lambda\to\pi_\lambda
    =\operatorname{Fun}\operatorname{Conn}(\Omega^{-\rho})^\lambda_D
\]
是单射，且其像含于 $\operatorname{Fun}\Op^{\check\lambda}_{{}^LG}\subset\operatorname{Fun}\operatorname{Conn}(\Omega^{-\rho})^\lambda_D$。如 §9.6.3 所说明（见图 (9.6-8) 之后的讨论），这证明了定理 9.6.1。

% [Tutor Note: p.304 的收尾把三条线索合在一起：(1) $\iota_\lambda$ 单射（定理 9.6.6）给出 (9.6-23) 单射；(2) 屏蔽算子 $\overline{S}^{(\lambda_i)}_i$ 的核给出像落在 $\operatorname{Fun}\Op^{\check\lambda}_{{}^LG}$ 内（命题 9.6.7 + 命题 9.6.3）；(3) §9.6.3 的图 (9.6-8) 给出反向的包含链。三点合起来即定理 9.6.1。注意屏蔽算子的 Fourier 指标是 $-(λ_i+1)$，与 (9.6-15) 的 $\widetilde{V}_i[\lambda_i+1]$ 指标一致，这是全节最后必须核对的一处。]





